\documentclass{article}
\usepackage[utf8]{inputenc}
\usepackage{amsmath, amssymb, graphicx, xcolor, fancyvrb}

% --- UBAH JUDUL, NAMA, DAN NIM ANDA DI SINI ---
\title{TugasAplikom8Adi25030630001}
\author{Adi Kurniawan \\ NIM: 25030630001}
\date{\today}

% --- PENERJEMAH OTOMATIS KODE EMT (JANGAN DIUBAH) ---
\newenvironment{eulernotebook}{\par}{\par}
\newcommand{\eulerheading}[1]{\section{#1}}
\newcommand{\eulersubheading}[1]{\subsection{#1}}
\newenvironment{eulercomment}{\par\medskip}{\par\medskip}
\newenvironment{eulerformula}{}{}
\DefineVerbatimEnvironment{eulerprompt}{Verbatim}{formatcom=\color{red}}
\DefineVerbatimEnvironment{euleroutput}{Verbatim}{formatcom=\color{blue}}
\DefineVerbatimEnvironment{eulerttcomment}{Verbatim}{formatcom=\color{teal}}
\newcommand{\eulerimg}[2]{\begin{center}\includegraphics[width=0.8\textwidth]{#2}\end{center}}

\begin{document}
\maketitle

% ====================================================================
% TEMPELKAN KODE DARI \begin{eulernotebook} SAMPAI \end{eulernotebook}
% DI BAWAH BARIS INI:
\begin{eulernotebook}
\eulerheading{EMT untuk Statistika}
\begin{eulercomment}
Di dalam notebook ini, kami mendemonstrasikan plot, uji, dan
distribusi statistik utama di Euler.

Mari kita mulai dengan beberapa statistik deskriptif. Ini bukanlah
pengantar statistika, sehingga Anda mungkin memerlukan sedikit latar
belakang pengetahuan untuk memahami detailnya.

Misalkan terdapat data pengukuran berikut. Kita ingin menghitung nilai
rata-rata (mean) dan simpangan baku terukur (measured standard
deviation).
\end{eulercomment}
\begin{eulerprompt}
>M=[1000,1004,998,997,1002,1001,998,1004,998,997]; ...
>median(M), mean(M), dev(M),
\end{eulerprompt}
\begin{euleroutput}
  999
  999.9
  2.72641400622
\end{euleroutput}
\begin{eulercomment}
Kita dapat membuat diagram kotak-garis (box-and-whiskers plot) untuk
data tersebut. Dalam kasus kita, tidak terdapat pencilan (outliers).
\end{eulercomment}
\begin{eulerprompt}
>aspect(1.75); boxplot(M):
\end{eulerprompt}
\eulerimg{15}{Image Use/TugasAplikom8Adi25030630001-001.png}
\begin{eulercomment}
Kita menghitung probabilitas suatu nilai lebih besar dari 1005, dengan
mengasumsikan nilai yang diukur berasal dari distribusi normal.

Semua fungsi untuk distribusi di Euler berakhiran ...dis dan
menghitung fungsi distribusi probabilitas kumulatif (cumulative
probability distribution / CPF).

\end{eulercomment}
\begin{eulerformula}
\[
\text{normaldis(x,m,d)}=\int_{-\infty}^x \frac{1}{d\sqrt{2\pi}}e^{-\frac{1}{2}(\frac{t-m}{d})^2}\ dt.
\]
\end{eulerformula}
\begin{eulercomment}
Kita mencetak hasilnya dalam \% dengan akurasi 2 digit menggunakan
fungsi print.
\end{eulercomment}
\begin{eulerprompt}
>print((1-normaldis(1005,mean(M),dev(M)))*100,2,unit=" %")
\end{eulerprompt}
\begin{euleroutput}
        3.07 %
\end{euleroutput}
\begin{eulercomment}
Untuk contoh berikutnya, kita mengasumsikan jumlah pria berikut pada
rentang ukuran yang diberikan.
\end{eulercomment}
\begin{eulerprompt}
>r=155.5:4:187.5; v=[22,71,136,169,139,71,32,8];
\end{eulerprompt}
\begin{eulercomment}
Berikut adalah plot dari distribusi tersebut.
\end{eulercomment}
\begin{eulerprompt}
>plot2d(r,v,a=150,b=200,c=0,d=190,bar=1,style="\(\backslash\)/"):
\end{eulerprompt}
\eulerimg{15}{Image Use/TugasAplikom8Adi25030630001-003.png}
\begin{eulercomment}
Kita dapat memasukkan data mentah tersebut ke dalam sebuah tabel.

Tabel adalah metode untuk menyimpan data statistik. Tabel kita harus
memuat tiga kolom: awal rentang, akhir rentang, dan jumlah pria dalam
rentang tersebut.

Tabel dapat dicetak dengan menyertakan tajuk kolom (*headers*). Kita
menggunakan vektor string untuk mengatur tajuk-tajuk tersebut.
\end{eulercomment}
\begin{eulerprompt}
>T:=r[1:8]' | r[2:9]' | v'; writetable(T,labc=["BB","BA","Frek"])
\end{eulerprompt}
\begin{euleroutput}
          BB        BA      Frek
       155.5     159.5        22
       159.5     163.5        71
       163.5     167.5       136
       167.5     171.5       169
       171.5     175.5       139
       175.5     179.5        71
       179.5     183.5        32
       183.5     187.5         8
\end{euleroutput}
\begin{eulercomment}
Jika kita memerlukan nilai rata-rata dan statistik ukuran lainnya,
kita perlu menghitung titik tengah dari rentang tersebut. Kita dapat
menggunakan dua kolom pertama dari tabel kita untuk hal ini.

Simbol "\textbar{}" digunakan untuk memisahkan kolom, fungsi "writetable"
digunakan untuk menulis tabel, dengan opsi "labc" untuk menentukan
tajuk kolom.
\end{eulercomment}
\begin{eulerprompt}
>(T[,1]+T[,2])/2 // the midpoint of each interval
\end{eulerprompt}
\begin{euleroutput}
          157.5 
          161.5 
          165.5 
          169.5 
          173.5 
          177.5 
          181.5 
          185.5 
\end{euleroutput}
\begin{eulercomment}
Namun, lebih mudah untuk melipat rentang tersebut dengan vektor
[1/2,1/2].
\end{eulercomment}
\begin{eulerprompt}
>M=fold(r,[0.5,0.5])
\end{eulerprompt}
\begin{euleroutput}
  [157.5,  161.5,  165.5,  169.5,  173.5,  177.5,  181.5,  185.5]
\end{euleroutput}
\begin{eulercomment}
Sekarang kita dapat menghitung rata-rata dan simpangan dari sampel
dengan frekuensi yang diberikan.
\end{eulercomment}
\begin{eulerprompt}
>\{m,d\}=meandev(M,v); m, d,
\end{eulerprompt}
\begin{euleroutput}
  169.901234568
  5.98912964449
\end{euleroutput}
\begin{eulercomment}
Mari kita tambahkan distribusi normal dari nilai-nilai tersebut ke
diagram batang di atas. Rumus untuk distribusi normal dengan rata-rata
\textdollar{}m\textdollar{} dan simpangan baku \textdollar{}d\textdollar{} adalah:

\end{eulercomment}
\begin{eulerformula}
\[
y=\frac{1}{d\sqrt{2\pi}}e^{\frac{-(x-m)^2}{2d^2}}.
\]
\end{eulerformula}
\begin{eulercomment}
Karena nilainya berada di antara 0 dan 1, untuk memplotnya pada
diagram batang, nilai tersebut harus dikalikan dengan 4 kali jumlah
total data.
\end{eulercomment}
\begin{eulerprompt}
> plot2d("qnormal(x,m,d)*sum(v)*4", ...
>  xmin=min(r),xmax=max(r),thickness=3,add=1):
\end{eulerprompt}
\eulerimg{15}{Image Use/TugasAplikom8Adi25030630001-005.png}
\eulerheading{Tables}
\begin{eulercomment}
Di direktori notebook ini, Anda akan menemukan sebuah file yang berisi
tabel. Data tersebut merepresentasikan hasil dari sebuah survei.
Berikut adalah empat baris pertama dari file tersebut. Data tersebut
berasal dari buku online berbahasa Jerman berjudul "Einführung in die
Statistik mit R" karya A. Handl.
\end{eulercomment}
\begin{eulerprompt}
>printfile("table.dat",4);
\end{eulerprompt}
\begin{euleroutput}
  Person Sex Age Titanic Evaluation Tip Problem
  1 m 30 n . 1.80 n
  2 f 23 y g 1.80 n
  3 f 26 y g 1.80 y
\end{euleroutput}
\begin{eulercomment}
Tabel tersebut berisi 7 kolom angka atau token (*string*). Kita ingin
membaca tabel dari file tersebut. Pertama, kita menggunakan terjemahan
kita sendiri untuk token-token tersebut.

Untuk itu, kita mendefinisikan himpunan token. Fungsi `strtokens()`
mengambil vektor *string* yang berisi token dari *string* yang
diberikan.
\end{eulercomment}
\begin{eulerprompt}
>mf:=["m","f"]; yn:=["y","n"]; ev:=strtokens("g vg m b vb");
\end{eulerprompt}
\begin{eulercomment}
Sekarang kita membaca tabel dengan terjemahan tersebut.

Argumen `tok2`, `tok4` dan seterusnya adalah terjemahan untuk
kolom-kolom pada tabel. Argumen-argumen ini tidak ada di dalam daftar
parameter `readtable()`, sehingga Anda perlu menyediakannya
menggunakan tanda `:=`.
\end{eulercomment}
\begin{eulerprompt}
>\{MT,hd\}=readtable("table.dat",tok2:=mf,tok4:=yn,tok5:=ev,tok7:=yn);
>load over statistics;
\end{eulerprompt}
\begin{eulercomment}
Untuk pencetakan, kita perlu menentukan himpunan token yang sama. Kita
hanya mencetak empat baris pertama.
\end{eulercomment}
\begin{eulerprompt}
>writetable(MT[1:10],labc=hd,wc=5,tok2:=mf,tok4:=yn,tok5:=ev,tok7:=yn);
\end{eulerprompt}
\begin{euleroutput}
   Person  Sex  Age Titanic Evaluation  Tip Problem
        1    m   30       n          .  1.8       n
        2    f   23       y          g  1.8       n
        3    f   26       y          g  1.8       y
        4    m   33       n          .  2.8       n
        5    m   37       n          .  1.8       n
        6    m   28       y          g  2.8       y
        7    f   31       y         vg  2.8       n
        8    m   23       n          .  0.8       n
        9    f   24       y         vg  1.8       y
       10    m   26       n          .  1.8       n
\end{euleroutput}
\begin{eulercomment}
Titik-titik "." merepresentasikan nilai yang tidak tersedia.

Jika kita tidak ingin menentukan token untuk terjemahan terlebih
dahulu, kita hanya perlu menentukan kolom mana yang berisi token dan
bukan angka.
\end{eulercomment}
\begin{eulerprompt}
>ctok=[2,4,5,7]; \{MT,hd,tok\}=readtable("table.dat",ctok=ctok);
\end{eulerprompt}
\begin{eulercomment}
Fungsi `readtable()` sekarang mengembalikan sekumpulan token.
\end{eulercomment}
\begin{eulerprompt}
>tok
\end{eulerprompt}
\begin{euleroutput}
  m
  n
  f
  y
  g
  vg
\end{euleroutput}
\begin{eulercomment}
Tabel tersebut berisi entri-entri dari file dengan token yang telah
diterjemahkan menjadi angka.

*String* khusus `NA="."` diinterpretasikan sebagai "Not Available",
dan menghasilkan nilai `NAN` (*not a number*) di dalam tabel.
Terjemahan ini dapat diubah menggunakan parameter `NA` dan `NAval`.
\end{eulercomment}
\begin{eulerprompt}
>MT[1]
\end{eulerprompt}
\begin{euleroutput}
  [1,  1,  30,  2,  NAN,  1.8,  2]
\end{euleroutput}
\begin{eulercomment}
Berikut adalah isi tabel dengan angka yang tidak diterjemahkan.
\end{eulercomment}
\begin{eulerprompt}
>writetable(MT,wc=5)
\end{eulerprompt}
\begin{euleroutput}
      1    1   30    2    .  1.8    2
      2    3   23    4    5  1.8    2
      3    3   26    4    5  1.8    4
      4    1   33    2    .  2.8    2
      5    1   37    2    .  1.8    2
      6    1   28    4    5  2.8    4
      7    3   31    4    6  2.8    2
      8    1   23    2    .  0.8    2
      9    3   24    4    6  1.8    4
     10    1   26    2    .  1.8    2
     11    3   23    4    6  1.8    4
     12    1   32    4    5  1.8    2
     13    1   29    4    6  1.8    4
     14    3   25    4    5  1.8    4
     15    3   31    4    5  0.8    2
     16    1   26    4    5  2.8    2
     17    1   37    2    .  3.8    2
     18    1   38    4    5    .    2
     19    3   29    2    .  3.8    2
     20    3   28    4    6  1.8    2
     21    3   28    4    1  2.8    4
     22    3   28    4    6  1.8    4
     23    3   38    4    5  2.8    2
     24    3   27    4    1  1.8    4
     25    1   27    2    .  2.8    4
\end{euleroutput}
\begin{eulercomment}
Demi kepraktisan, Anda dapat memasukkan keluaran dari readtable() ke
dalam sebuah list.
\end{eulercomment}
\begin{eulerprompt}
>Table=\{\{readtable("table.dat",ctok=ctok)\}\};
\end{eulerprompt}
\begin{eulercomment}
Dengan menggunakan kolom token yang sama dan token yang dibaca dari
file, kita dapat mencetak tabel tersebut. Kita bisa menentukan ctok,
tok, dan sebagainya, atau menggunakan list Table.
\end{eulercomment}
\begin{eulerprompt}
>writetable(Table,ctok=ctok,wc=5);
\end{eulerprompt}
\begin{euleroutput}
   Person  Sex  Age Titanic Evaluation  Tip Problem
        1    m   30       n          .  1.8       n
        2    f   23       y          g  1.8       n
        3    f   26       y          g  1.8       y
        4    m   33       n          .  2.8       n
        5    m   37       n          .  1.8       n
        6    m   28       y          g  2.8       y
        7    f   31       y         vg  2.8       n
        8    m   23       n          .  0.8       n
        9    f   24       y         vg  1.8       y
       10    m   26       n          .  1.8       n
       11    f   23       y         vg  1.8       y
       12    m   32       y          g  1.8       n
       13    m   29       y         vg  1.8       y
       14    f   25       y          g  1.8       y
       15    f   31       y          g  0.8       n
       16    m   26       y          g  2.8       n
       17    m   37       n          .  3.8       n
       18    m   38       y          g    .       n
       19    f   29       n          .  3.8       n
       20    f   28       y         vg  1.8       n
       21    f   28       y          m  2.8       y
       22    f   28       y         vg  1.8       y
       23    f   38       y          g  2.8       n
       24    f   27       y          m  1.8       y
       25    m   27       n          .  2.8       y
\end{euleroutput}
\begin{eulercomment}
Fungsi tablecol() mengembalikan nilai-nilai dari kolom tabel, dengan
melewatkan baris apa pun yang memiliki nilai NAN ("." di dalam file),
serta indeks dari kolom-kolom yang memuat nilai-nilai tersebut.
\end{eulercomment}
\begin{eulerprompt}
>\{c,i\}=tablecol(MT,[5,6]);
\end{eulerprompt}
\begin{eulercomment}
Kita dapat menggunakan ini untuk mengekstrak kolom dari tabel guna
membuat tabel baru.
\end{eulercomment}
\begin{eulerprompt}
>j=[1,5,6]; writetable(MT[i,j],labc=hd[j],ctok=[2],tok=tok)
\end{eulerprompt}
\begin{euleroutput}
      Person Evaluation       Tip
           2          g       1.8
           3          g       1.8
           6          g       2.8
           7         vg       2.8
           9         vg       1.8
          11         vg       1.8
          12          g       1.8
          13         vg       1.8
          14          g       1.8
          15          g       0.8
          16          g       2.8
          20         vg       1.8
          21          m       2.8
          22         vg       1.8
          23          g       2.8
          24          m       1.8
\end{euleroutput}
\begin{eulercomment}
Tentu saja, dalam kasus ini kita perlu mengekstrak tabel itu sendiri
dari list Table.
\end{eulercomment}
\begin{eulerprompt}
>MT=Table[1];
\end{eulerprompt}
\begin{eulercomment}
Tentu saja, kita juga dapat menggunakannya untuk menentukan nilai
rata-rata dari suatu kolom atau nilai statistik lainnya.
\end{eulercomment}
\begin{eulerprompt}
>mean(tablecol(MT,6))
\end{eulerprompt}
\begin{euleroutput}
  2.175
\end{euleroutput}
\begin{eulercomment}
Fungsi getstatistics() mengembalikan elemen-elemen dalam sebuah vektor
beserta jumlah kemunculannya. Kita menerapkannya pada nilai "m" dan
"f" di kolom kedua pada tabel kita.
\end{eulercomment}
\begin{eulerprompt}
>\{xu,count\}=getstatistics(tablecol(MT,2)); xu, count,
\end{eulerprompt}
\begin{euleroutput}
  [1,  3]
  [12,  13]
\end{euleroutput}
\begin{eulercomment}
Kita dapat mencetak hasilnya ke dalam sebuah tabel baru.
\end{eulercomment}
\begin{eulerprompt}
>writetable(count',labr=tok[xu])
\end{eulerprompt}
\begin{euleroutput}
           m        12
           f        13
\end{euleroutput}
\begin{eulercomment}
Fungsi selecttable() mengembalikan tabel baru dengan nilai-nilai pada
satu kolom yang dipilih dari sebuah vektor indeks. Pertama, kita
mencari indeks dari dua nilai kita di dalam tabel token.
\end{eulercomment}
\begin{eulerprompt}
>v:=indexof(tok,["g","vg"])
\end{eulerprompt}
\begin{euleroutput}
  [5,  6]
\end{euleroutput}
\begin{eulercomment}
Sekarang kita dapat memilih baris-baris tabel yang memiliki salah satu
nilai di dalam v pada baris ke-5 mereka.
\end{eulercomment}
\begin{eulerprompt}
>MT1:=MT[selectrows(MT,5,v)]; i:=sortedrows(MT1,5);
\end{eulerprompt}
\begin{eulercomment}
Sekarang kita dapat mencetak tabel tersebut, dengan nilai-nilai yang
telah diekstrak dan diurutkan pada kolom ke-5.
\end{eulercomment}
\begin{eulerprompt}
>writetable(MT1[i],labc=hd,ctok=ctok,tok=tok,wc=7);
\end{eulerprompt}
\begin{euleroutput}
   Person    Sex    Age Titanic Evaluation    Tip Problem
        2      f     23       y          g    1.8       n
        3      f     26       y          g    1.8       y
        6      m     28       y          g    2.8       y
       18      m     38       y          g      .       n
       16      m     26       y          g    2.8       n
       15      f     31       y          g    0.8       n
       12      m     32       y          g    1.8       n
       23      f     38       y          g    2.8       n
       14      f     25       y          g    1.8       y
        9      f     24       y         vg    1.8       y
        7      f     31       y         vg    2.8       n
       20      f     28       y         vg    1.8       n
       22      f     28       y         vg    1.8       y
       13      m     29       y         vg    1.8       y
       11      f     23       y         vg    1.8       y
\end{euleroutput}
\begin{eulercomment}
Untuk statistik berikutnya, kita ingin menghubungkan dua kolom dari
tabel tersebut. Jadi, kita mengekstrak kolom 2 dan 4 lalu mengurutkan
tabelnya.
\end{eulercomment}
\begin{eulerprompt}
>i=sortedrows(MT,[2,4]);  ...
>  writetable(tablecol(MT[i],[2,4])',ctok=[1,2],tok=tok)
\end{eulerprompt}
\begin{euleroutput}
           m         n
           m         n
           m         n
           m         n
           m         n
           m         n
           m         n
           m         y
           m         y
           m         y
           m         y
           m         y
           f         n
           f         y
           f         y
           f         y
           f         y
           f         y
           f         y
           f         y
           f         y
           f         y
           f         y
           f         y
           f         y
\end{euleroutput}
\begin{eulercomment}
Dengan getstatistics(), kita juga dapat menghubungkan jumlah
kemunculan pada dua kolom tabel satu sama lain.
\end{eulercomment}
\begin{eulerprompt}
>MT24=tablecol(MT,[2,4]); ...
>\{xu1,xu2,count\}=getstatistics(MT24[1],MT24[2]); ...
>writetable(count,labr=tok[xu1],labc=tok[xu2])
\end{eulerprompt}
\begin{euleroutput}
                     n         y
           m         7         5
           f         1        12
\end{euleroutput}
\begin{eulercomment}
Sebuah tabel dapat ditulis ke dalam sebuah file.
\end{eulercomment}
\begin{eulerprompt}
>filename="test.dat"; ...
>writetable(count,labr=tok[xu1],labc=tok[xu2],file=filename);
\end{eulerprompt}
\begin{eulercomment}
Kemudian kita dapat membaca tabel tersebut dari file.
\end{eulercomment}
\begin{eulerprompt}
>\{MT2,hd,tok2,hdr\}=readtable(filename,>clabs,>rlabs); ...
>writetable(MT2,labr=hdr,labc=hd)
\end{eulerprompt}
\begin{euleroutput}
                     n         y
           m         7         5
           f         1        12
\end{euleroutput}
\begin{eulercomment}
Dan menghapus file tersebut.
\end{eulercomment}
\begin{eulerprompt}
>fileremove(filename);
\end{eulerprompt}
\eulerheading{Distributions}
\begin{eulercomment}
Dengan plot2d, ada cara yang sangat mudah untuk memplot distribusi
data eksperimen.
\end{eulercomment}
\begin{eulerprompt}
>p=normal(1,1000); //1000 random normal-distributed sample p
>plot2d(p,distribution=20,style="\(\backslash\)/"); // plot the random sample p
>plot2d("qnormal(x,0,1)",add=1): // add the standard normal distribution plot
\end{eulerprompt}
\eulerimg{27}{Image Use/TugasAplikom8Adi25030630001-006.png}
\begin{eulercomment}
Harap perhatikan perbedaan antara plot batang (sampel) dan kurva
normal (distribusi sebenarnya). Masukkan kembali ketiga perintah
tersebut untuk melihat hasil sampling lainnya.
\end{eulercomment}
\begin{eulercomment}
Berikut adalah perbandingan dari 10 simulasi yang masing-masing
terdiri dari 1000 nilai berdistribusi normal menggunakan yang disebut
*box plot* (diagram kotak garis). Plot ini menunjukkan median, kuartil
25\% dan 75\%, nilai minimum dan maksimum, serta pencilan (*outlier*).
\end{eulercomment}
\begin{eulerprompt}
>p=normal(10,1000); boxplot(p):
\end{eulerprompt}
\eulerimg{27}{Image Use/TugasAplikom8Adi25030630001-007.png}
\begin{eulercomment}
Untuk menghasilkan bilangan bulat acak, Euler memiliki `intrandom`.
Mari kita simulasikan lemparan dadu dan plot distribusinya.

Kita menggunakan fungsi `getmultiplicities(v,x)`, yang menghitung
seberapa sering elemen-elemen pada `v` muncul di dalam `x`. Kemudian
kita memplot hasilnya menggunakan `columnsplot()`.
\end{eulercomment}
\begin{eulerprompt}
>k=intrandom(1,6000,6);  ...
>columnsplot(getmultiplicities(1:6,k));  ...
>ygrid(1000,color=red):
\end{eulerprompt}
\eulerimg{27}{Image Use/TugasAplikom8Adi25030630001-008.png}
\begin{eulercomment}
Meskipun `intrandom(n,m,k)` menghasilkan bilangan bulat terdistribusi
seragam dari 1 sampai `k`, kita tetap bisa menggunakan distribusi
bilangan bulat lainnya dengan `randpint()`.

Pada contoh berikut, probabilitas untuk 1, 2, dan 3 masing-masing
adalah 0,4, 0,1, dan 0,5.
\end{eulercomment}
\begin{eulerprompt}
>randpint(1,1000,[0.4,0.1,0.5]); getmultiplicities(1:3,%)
\end{eulerprompt}
\begin{euleroutput}
  [394,  107,  499]
\end{euleroutput}
\begin{eulercomment}
Euler dapat menghasilkan nilai acak dari lebih banyak distribusi.
Silakan lihat ke dalam referensi.

Sebagai contoh, kita mencoba distribusi eksponensial. Suatu peubah
acak kontinu \textdollar{}X\textdollar{} dikatakan memiliki distribusi eksponensial jika
fungsi kepekatan peluangnya (PDF) diberikan oleh\\
\end{eulercomment}
\begin{eulerformula}
\[
f_X(x)=\lambda e^{-\lambda x},\quad x>0,\quad \lambda>0,
\]
\end{eulerformula}
\begin{eulercomment}
dengan parameter\\
\end{eulercomment}
\begin{eulerformula}
\[
\lambda=\frac{1}{\mu},\quad \mu \text{ is the mean, and denoted by } X \sim \text{Exponential}(\lambda).
\]
\end{eulerformula}
\begin{eulerprompt}
>plot2d(randexponential(1,1000,2),>distribution):
\end{eulerprompt}
\eulerimg{15}{Image Use/TugasAplikom8Adi25030630001-011.png}
\begin{eulercomment}
Untuk banyak distribusi, Euler dapat menghitung fungsi distribusi dan
inversnya.
\end{eulercomment}
\begin{eulerprompt}
>plot2d("normaldis",-4,4): 
\end{eulerprompt}
\eulerimg{27}{Image Use/TugasAplikom8Adi25030630001-012.png}
\begin{eulercomment}
Berikut ini adalah salah satu cara untuk memplot kuantil.
\end{eulercomment}
\begin{eulerprompt}
>plot2d("qnormal(x,1,1.5)",-4,6);  ...
>plot2d("qnormal(x,1,1.5)",a=2,b=5,>add,>filled):
\end{eulerprompt}
\eulerimg{27}{Image Use/TugasAplikom8Adi25030630001-013.png}
\begin{eulerformula}
\[
\text{normaldis(x,m,d)}=\int_{-\infty}^x \frac{1}{d\sqrt{2\pi}}e^{-\frac{1}{2}(\frac{t-m}{d})^2}\ dt.
\]
\end{eulerformula}
\begin{eulercomment}
Probabilitas berada di area hijau adalah sebagai berikut.
\end{eulercomment}
\begin{eulerprompt}
>normaldis(5,1,1.5)-normaldis(2,1,1.5)
\end{eulerprompt}
\begin{euleroutput}
  0.248662156979
\end{euleroutput}
\begin{eulercomment}
Hal ini dapat dihitung secara numerik dengan integral berikut.\\
\end{eulercomment}
\begin{eulerformula}
\[
\int_2^5 \frac{1}{1.5\sqrt{2\pi}}e^{-\frac{1}{2}(\frac{x-1}{1.5})^2}\ dx.
\]
\end{eulerformula}
\begin{eulerprompt}
>gauss("qnormal(x,1,1.5)",2,5)
\end{eulerprompt}
\begin{euleroutput}
  0.248662156979
\end{euleroutput}
\begin{eulercomment}
Mari kita bandingkan distribusi binomial dengan distribusi normal yang
memiliki rataan (mean) dan simpangan baku yang sama. Fungsi
invbindis() menyelesaikan interpolasi linear di antara nilai-nilai
bilangan bulat.
\end{eulercomment}
\begin{eulerprompt}
>invbindis(0.95,1000,0.5), invnormaldis(0.95,500,0.5*sqrt(1000))
\end{eulerprompt}
\begin{euleroutput}
  525.516721219
  526.007419394
\end{euleroutput}
\begin{eulercomment}
Fungsi qdis() adalah fungsi kepekatan (densitas) dari distribusi
khi-kuadrat (chi-square). Seperti biasa, Euler memetakan vektor ke
fungsi ini. Dengan demikian, kita memperoleh plot dari semua
distribusi khi-kuadrat dengan derajat kebebasan 5 sampai 30 secara
mudah dengan cara berikut.
\end{eulercomment}
\begin{eulerprompt}
>plot2d("qchidis(x,(5:5:50)')",0,50):
\end{eulerprompt}
\eulerimg{27}{Image Use/TugasAplikom8Adi25030630001-016.png}
\begin{eulercomment}
Euler memiliki fungsi-fungsi yang akurat untuk mengevaluasi
distribusi. Mari kita periksa chidis() dengan sebuah integral.

Penamaannya diupayakan agar konsisten. Misalnya:

- distribusi khi-kuadrat (chi-square) adalah chidis(),\\
- ungsi inversnya adalah invchidis(),\\
- fungsi kepekatannya (densitas) adalah qchidis().

Komplemen dari distribusi tersebut (ekor atas / upper tail) adalah
chicdis().
\end{eulercomment}
\begin{eulerprompt}
>chidis(1.5,2), integrate("qchidis(x,2)",0,1.5)
\end{eulerprompt}
\begin{euleroutput}
  0.527633447259
  0.527633447259
\end{euleroutput}
\eulerheading{Discrete Distributions}
\begin{eulercomment}
Untuk mendefinisikan distribusi diskret Anda sendiri, Anda dapat
menggunakan metode berikut.

Pertama, kita menetapkan fungsi distribusinya.
\end{eulercomment}
\begin{eulerprompt}
>wd = 0|((1:6)+[-0.01,0.01,0,0,0,0])/6
\end{eulerprompt}
\begin{euleroutput}
  [0,  0.165,  0.335,  0.5,  0.666667,  0.833333,  1]
\end{euleroutput}
\begin{eulercomment}
Artinya adalah bahwa dengan probabilitas wd[i+1]-wd[i], kita
menghasilkan nilai acak i.

Hal ini hampir menyerupai distribusi seragam (uniform distribution).
Mari kita definisikan pembangkit bilangan acak untuk ini. Fungsi
find(v,x) mencari nilai x di dalam vektor v. Fungsi ini juga berlaku
untuk vektor x.
\end{eulercomment}
\begin{eulerprompt}
>function wrongdice (n,m) := find(wd,random(n,m))
\end{eulerprompt}
\begin{eulercomment}
Galatnya sangat halus sehingga kita baru melihatnya dengan iterasi
yang sangat banyak.
\end{eulercomment}
\begin{eulerprompt}
>columnsplot(getmultiplicities(1:6,wrongdice(1,1000000))):
\end{eulerprompt}
\eulerimg{27}{Image Use/TugasAplikom8Adi25030630001-017.png}
\begin{eulercomment}
Berikut adalah fungsi sederhana untuk memeriksa distribusi seragam
dari nilai-nilai 1...K di dalam v. Kita menerima hasilnya jika untuk
semua frekuensi

\end{eulercomment}
\begin{eulerformula}
\[
\left|f_i-\frac{1}{K}\right| < \frac{\delta}{\sqrt{n}}.
\]
\end{eulerformula}
\begin{eulerprompt}
>function checkrandom (v, delta=1) ...
\end{eulerprompt}
\begin{eulerudf}
    K=max(v); n=cols(v);
    fr=getfrequencies(v,1:K);
    return max(fr/n-1/K)<delta/sqrt(n);
    endfunction
\end{eulerudf}
\begin{eulercomment}
Memang benar bahwa fungsi tersebut menolak distribusi seragam.
\end{eulercomment}
\begin{eulerprompt}
>checkrandom(wrongdice(1,1000000))
\end{eulerprompt}
\begin{euleroutput}
  0
\end{euleroutput}
\begin{eulercomment}
Dan fungsi tersebut menerima pembangkit bilangan acak bawaan.
\end{eulercomment}
\begin{eulerprompt}
>checkrandom(intrandom(1,1000000,6))
\end{eulerprompt}
\begin{euleroutput}
  1
\end{euleroutput}
\begin{eulercomment}
Kita dapat menghitung distribusi binomial. Pertama, terdapat fungsi
binomialsum(), yang mengembalikan probabilitas dari i atau kurang
keberhasilan dari n percobaan.
\end{eulercomment}
\begin{eulerprompt}
>bindis(410,1000,0.4)
\end{eulerprompt}
\begin{euleroutput}
  0.751401349654
\end{euleroutput}
\begin{eulercomment}
Fungsi invers Beta digunakan untuk menghitung selang kepercayaan
Clopper-Pearson bagi parameter p. Tingkat signifikansi (atau taraf)
bawaannya adalah alpha.

Makna dari selang ini adalah jika p berada di luar selang tersebut,
maka hasil yang teramati sebesar 410 dari 1000 merupakan kejadian yang
jarang.
\end{eulercomment}
\begin{eulerprompt}
>clopperpearson(410,1000)
\end{eulerprompt}
\begin{euleroutput}
  [0.37932,  0.441212]
\end{euleroutput}
\begin{eulercomment}
Perintah-perintah berikut merupakan cara langsung untuk memperoleh
hasil di atas. Namun, untuk n yang besar, penjumlahan langsung tidak
akurat dan lambat.
\end{eulercomment}
\begin{eulerprompt}
>p=0.4; i=0:410; n=1000; sum(bin(n,i)*p^i*(1-p)^(n-i))
\end{eulerprompt}
\begin{euleroutput}
  0.751401349655
\end{euleroutput}
\begin{eulercomment}
Omong-omong, invbinsum() menghitung invers dari binomialsum().
\end{eulercomment}
\begin{eulerprompt}
>invbindis(0.75,1000,0.4)
\end{eulerprompt}
\begin{euleroutput}
  409.932733047
\end{euleroutput}
\begin{eulercomment}
Dalam permainan Bridge, kita mengasumsikan terdapat 5 kartu yang
tersisa (dari total 52 kartu) di dua tangan pemain (26 kartu). Mari
kita hitung probabilitas terjadinya distribusi yang lebih buruk
daripada 3:2 (misalnya 0:5, 1:4, 4:1, atau 5:0).
\end{eulercomment}
\begin{eulerprompt}
>2*hypergeomsum(1,5,13,26)
\end{eulerprompt}
\begin{euleroutput}
  0.321739130435
\end{euleroutput}
\begin{eulercomment}
Terdapat juga simulasi untuk distribusi multinomial.
\end{eulercomment}
\begin{eulerprompt}
>randmultinomial(10,1000,[0.4,0.1,0.5])
\end{eulerprompt}
\begin{euleroutput}
            394            82           524 
            418            92           490 
            445            90           465 
            403           114           483 
            405            95           500 
            384           107           509 
            414            93           493 
            419            90           491 
            394           101           505 
            382           103           515 
\end{euleroutput}
\eulerheading{Plotting Data}
\begin{eulercomment}
Untuk memplot data, kita mencoba menggunakan hasil pemilu Jerman sejak
tahun 1990, yang diukur dalam jumlah kursi.
\end{eulercomment}
\begin{eulerprompt}
>BW := [ ...
>1990,662,319,239,79,8,17; ...
>1994,672,294,252,47,49,30; ...
>1998,669,245,298,43,47,36; ...
>2002,603,248,251,47,55,2; ...
>2005,614,226,222,61,51,54; ...
>2009,622,239,146,93,68,76; ...
>2013,631,311,193,0,63,64];
\end{eulerprompt}
\begin{eulercomment}
Untuk nama-nama partai, kita menggunakan sebuah string nama.
\end{eulercomment}
\begin{eulerprompt}
>P:=["CDU/CSU","SPD","FDP","Gr","Li"];
\end{eulerprompt}
\begin{eulercomment}
Mari kita cetak persentasenya dengan rapi.

Pertama, kita mengekstrak kolom-kolom yang diperlukan. Kolom 3 hingga
7 adalah perolehan kursi masing-masing partai, dan kolom 2 adalah
jumlah total kursi. Kolom 1 merupakan tahun pemilu.
\end{eulercomment}
\begin{eulerprompt}
>BT:=BW[,3:7]; BT:=BT/sum(BT); YT:=BW[,1]';
\end{eulerprompt}
\begin{eulercomment}
Kemudian kita mencetak statistik tersebut dalam bentuk tabel. Kita
menggunakan nama-nama (partai) sebagai tajuk kolom, dan tahun-tahun
sebagai tajuk baris. Lebar bawaan untuk kolom adalah wc=10, namun kita
lebih memilih keluaran yang lebih rapat. Kolom-kolom tersebut akan
diperlebar untuk menyesuaikan label kolom jika diperlukan.
\end{eulercomment}
\begin{eulerprompt}
>writetable(BT*100,wc=6,dc=0,>fixed,labc=P,labr=YT)
\end{eulerprompt}
\begin{euleroutput}
         CDU/CSU   SPD   FDP    Gr    Li
    1990      48    36    12     1     3
    1994      44    38     7     7     4
    1998      37    45     6     7     5
    2002      41    42     8     9     0
    2005      37    36    10     8     9
    2009      38    23    15    11    12
    2013      49    31     0    10    10
\end{euleroutput}
\begin{eulercomment}
Perkalian matriks berikut mengekstrak jumlah persentase dari dua
partai besar, yang menunjukkan bahwa partai-partai kecil telah
memperoleh pijakan/dukungan di parlemen hingga tahun 2009.
\end{eulercomment}
\begin{eulerprompt}
>BT1:=(BT.[1;1;0;0;0])'*100
\end{eulerprompt}
\begin{euleroutput}
  [84.29,  81.25,  81.1659,  82.7529,  72.9642,  61.8971,  79.8732]
\end{euleroutput}
\begin{eulercomment}
Terdapat juga plot statistik sederhana. Kita menggunakannya untuk
menampilkan garis dan titik secara bersamaan. Alternatifnya adalah
memanggil plot2d dua kali dengan argumen \textgreater{}add.
\end{eulercomment}
\begin{eulerprompt}
>statplot(YT,BT1,"b"):
\end{eulerprompt}
\eulerimg{27}{Image Use/TugasAplikom8Adi25030630001-019.png}
\begin{eulercomment}
Definisikan beberapa warna untuk masing-masing partai.
\end{eulercomment}
\begin{eulerprompt}
>CP:=[rgb(0.5,0.5,0.5),red,yellow,green,rgb(0.8,0,0)];
\end{eulerprompt}
\begin{eulercomment}
Sekarang kita dapat memplot hasil pemilu tahun 2009 dan perubahannya
ke dalam satu plot menggunakan figure. Kita dapat menambahkan sebuah
vektor kolom ke masing-masing plot.
\end{eulercomment}
\begin{eulerprompt}
>figure(2,1);  ...
>figure(1); columnsplot(BW[6,3:7],P,color=CP); ...
>figure(2); columnsplot(BW[6,3:7]-BW[5,3:7],P,color=CP);  ...
>figure(0):
\end{eulerprompt}
\eulerimg{27}{Image Use/TugasAplikom8Adi25030630001-020.png}
\begin{eulercomment}
Plot data menggabungkan baris-baris dari data statistik ke dalam satu
plot.
\end{eulercomment}
\begin{eulerprompt}
>J:=BW[,1]'; DP:=BW[,3:7]'; ...
>dataplot(YT,BT',color=CP);  ...
>labelbox(P,colors=CP,styles="[]",>points,w=0.2,x=0.3,y=0.4):
\end{eulerprompt}
\eulerimg{27}{Image Use/TugasAplikom8Adi25030630001-021.png}
\begin{eulercomment}
Plot kolom 3D menampilkan baris-baris dari data statistik dalam bentuk
kolom. Kita menyediakan label untuk baris-baris dan kolom-kolom
tersebut. angle adalah sudut pandang.
\end{eulercomment}
\begin{eulerprompt}
>columnsplot3d(BT,scols=P,srows=YT, ...
>  angle=30°,ccols=CP):
\end{eulerprompt}
\eulerimg{27}{Image Use/TugasAplikom8Adi25030630001-022.png}
\begin{eulercomment}
Representasi lainnya adalah mosaic plot. Perlu dicatat bahwa
kolom-kolom pada plot ini merepresentasikan kolom-kolom dari matriks.
Karena panjangnya label CDU/CSU, kita menggunakan jendela (window)
yang lebih kecil daripada biasanya.
\end{eulercomment}
\begin{eulerprompt}
>shrinkwindow(>smaller);  ...
>mosaicplot(BT',srows=YT,scols=P,color=CP,style="#"); ...
>shrinkwindow():
\end{eulerprompt}
\eulerimg{27}{Image Use/TugasAplikom8Adi25030630001-023.png}
\begin{eulercomment}
Kita juga dapat membuat diagram lingkaran (pie chart). Karena warna
hitam dan kuning membentuk sebuah koalisi, kita menyusun ulang urutan
elemen-elemennya.
\end{eulercomment}
\begin{eulerprompt}
>i=[1,3,5,4,2]; piechart(BW[6,3:7][i],color=CP[i],lab=P[i]):
\end{eulerprompt}
\eulerimg{27}{Image Use/TugasAplikom8Adi25030630001-024.png}
\begin{eulercomment}
Berikut adalah jenis plot lainnya.
\end{eulercomment}
\begin{eulerprompt}
>starplot(normal(1,10)+4,lab=1:10,>rays):
\end{eulerprompt}
\eulerimg{27}{Image Use/TugasAplikom8Adi25030630001-025.png}
\begin{eulercomment}
Beberapa plot pada plot2d sangat baik untuk statistika. Berikut adalah
impulse plot dari data acak yang berdistribusi seragam pada interval
[0,1].
\end{eulercomment}
\begin{eulerprompt}
>plot2d(makeimpulse(1:10,random(1,10)),>bar):
\end{eulerprompt}
\eulerimg{27}{Image Use/TugasAplikom8Adi25030630001-026.png}
\begin{eulercomment}
Namun untuk data yang berdistribusi eksponensial, kita mungkin
memerlukan plot logaritmik.
\end{eulercomment}
\begin{eulerprompt}
>logimpulseplot(1:10,-log(random(1,10))*10):
\end{eulerprompt}
\eulerimg{27}{Image Use/TugasAplikom8Adi25030630001-027.png}
\begin{eulercomment}
Fungsi columnsplot() lebih mudah digunakan karena hanya memerlukan
sebuah vektor nilai. Selain itu, fungsi ini dapat mengatur labelnya
sesuai dengan yang kita inginkan; hal ini telah kita demonstrasikan
sebelumnya dalam tutorial ini.

Berikut adalah aplikasi lainnya, di mana kita menghitung jumlah
karakter dalam sebuah kalimat dan memplot statistiknya.
\end{eulercomment}
\begin{eulerprompt}
>v=strtochar("the quick brown fox jumps over the lazy dog"); ...
>w=ascii("a"):ascii("z"); x=getmultiplicities(w,v); ...
>cw=[]; for k=w; cw=cw|char(k); end; ...
>columnsplot(x,lab=cw,width=0.05):
\end{eulerprompt}
\eulerimg{27}{Image Use/TugasAplikom8Adi25030630001-028.png}
\begin{eulercomment}
Kita juga dapat mengatur sumbu secara manual.
\end{eulercomment}
\begin{eulerprompt}
>n=10; p=0.4; i=0:n; x=bin(n,i)*p^i*(1-p)^(n-i); ...
>columnsplot(x,lab=i,width=0.05,<frame,<grid); ...
>yaxis(0,0:0.1:1,style="->",>left); xaxis(0,style="."); ...
>label("p",0,0.25), label("i",11,0); ...
>textbox(["Binomial distribution","with p=0.4"]):
\end{eulerprompt}
\eulerimg{27}{Image Use/TugasAplikom8Adi25030630001-029.png}
\begin{eulercomment}
Berikut ini adalah salah satu cara untuk memplot frekuensi kemunculan
angka-angka dalam sebuah vektor.

Kita membuat sebuah vektor yang berisi bilangan acak bulat dari 1
hingga 6.
\end{eulercomment}
\begin{eulerprompt}
>v:=intrandom(1,10,10)
\end{eulerprompt}
\begin{euleroutput}
  [5,  1,  8,  6,  6,  10,  9,  6,  8,  3]
\end{euleroutput}
\begin{eulercomment}
Kemudian ekstrak bilangan-bilangan unik di dalam v.
\end{eulercomment}
\begin{eulerprompt}
>vu:=unique(v)
\end{eulerprompt}
\begin{euleroutput}
  [1,  3,  5,  6,  8,  9,  10]
\end{euleroutput}
\begin{eulercomment}
Dan plot frekuensi tersebut ke dalam sebuah plot kolom.
\end{eulercomment}
\begin{eulerprompt}
>columnsplot(getmultiplicities(vu,v),lab=vu,style="/"):
\end{eulerprompt}
\eulerimg{27}{Image Use/TugasAplikom8Adi25030630001-030.png}
\begin{eulercomment}
Kita ingin mendemonstrasikan fungsi-fungsi untuk distribusi empiris
dari suatu nilai.
\end{eulercomment}
\begin{eulerprompt}
>x=normal(1,20);
\end{eulerprompt}
\begin{eulercomment}
Fungsi empdist(x,vs) memerlukan sebuah larik nilai yang sudah terurut.
Oleh karena itu, kita harus mengurutkan x terlebih dahulu sebelum
dapat menggunakannya.
\end{eulercomment}
\begin{eulerprompt}
>xs=sort(x);
\end{eulerprompt}
\begin{eulercomment}
Kemudian kita memplot distribusi empiris dan beberapa batang kerapatan
ke dalam satu plot. Alih-alih menggunakan bar plot untuk
distribusinya, kali ini kita menggunakan sawtooth plot.
\end{eulercomment}
\begin{eulerprompt}
>figure(2,1); ...
>figure(1); plot2d("empdist",-4,4;xs); ...
>figure(2); plot2d(histo(x,v=-4:0.2:4,<bar));  ...
>figure(0):
\end{eulerprompt}
\eulerimg{27}{Image Use/TugasAplikom8Adi25030630001-031.png}
\begin{eulercomment}
Plot sebar (scatter plot) mudah dibuat di Euler menggunakan point plot
yang biasa digunakan. Grafik berikut menunjukkan bahwa X dan X+Y jelas
berkorelasi positif.
\end{eulercomment}
\begin{eulerprompt}
>x=normal(1,100); plot2d(x,x+rotright(x),>points,style=".."):
\end{eulerprompt}
\eulerimg{27}{Image Use/TugasAplikom8Adi25030630001-032.png}
\begin{eulercomment}
Sering kali, kita ingin membandingkan dua sampel dari distribusi yang
berbeda. Hal ini dapat dilakukan dengan quantile-quantile-plot.

Untuk pengujian, kita mencoba distribusi Student-t dan distribusi
eksponensial.
\end{eulercomment}
\begin{eulerprompt}
>x=randt(1,1000,5); y=randnormal(1,1000,mean(x),dev(x)); ...
>plot2d("x",r=6,style="--",yl="normal",xl="student-t",>vertical); ...
>plot2d(sort(x),sort(y),>points,color=red,style="x",>add):
\end{eulerprompt}
\eulerimg{27}{Image Use/TugasAplikom8Adi25030630001-033.png}
\begin{eulercomment}
Plot tersebut dengan jelas menunjukkan bahwa nilai-nilai yang
berdistribusi normal cenderung lebih kecil pada ujung-ujung
ekstremnya.

Jika kita memiliki dua distribusi dengan ukuran yang berbeda, kita
dapat memperluas yang lebih kecil atau memperkecil yang lebih besar.
Fungsi berikut sangat baik untuk keduanya. Fungsi ini mengambil
nilai-nilai median dengan persentase antara 0 dan 1.
\end{eulercomment}
\begin{eulerprompt}
>function medianexpand (x,n) := median(x,p=linspace(0,1,n-1));
\end{eulerprompt}
\begin{eulercomment}
Mari kita bandingkan dua distribusi yang sama.
\end{eulercomment}
\begin{eulerprompt}
>x=random(1000); y=random(400); ...
>plot2d("x",0,1,style="--"); ...
>plot2d(sort(medianexpand(x,400)),sort(y),>points,color=red,style="x",>add):
\end{eulerprompt}
\eulerimg{27}{Image Use/TugasAplikom8Adi25030630001-034.png}
\eulerheading{Regression and Correlation}
\begin{eulercomment}
Regresi linear dapat dilakukan dengan fungsi polyfit() atau berbagai
fungsi fit lainnya.

Sebagai permulaan, kita menentukan garis regresi untuk data univariat
dengan polyfit(x,y,1).
\end{eulercomment}
\begin{eulerprompt}
>x=1:10; y=[2,3,1,5,6,3,7,8,9,8]; writetable(x'|y',labc=["x","y"])
\end{eulerprompt}
\begin{euleroutput}
           x         y
           1         2
           2         3
           3         1
           4         5
           5         6
           6         3
           7         7
           8         8
           9         9
          10         8
\end{euleroutput}
\begin{eulercomment}
Kita ingin membandingkan fit tak berbobot dan fit berbobot. Pertama,
koefisien-koefisien dari fit linear.
\end{eulercomment}
\begin{eulerprompt}
>p=polyfit(x,y,1)
\end{eulerprompt}
\begin{euleroutput}
  [0.733333,  0.812121]
\end{euleroutput}
\begin{eulercomment}
Sekarang, koefisien-koefisien dengan bobot yang menekankan nilai-nilai
terakhir.
\end{eulercomment}
\begin{eulerprompt}
>w &= "exp(-(x-10)^2/10)"; pw=polyfit(x,y,1,w=w(x))
\end{eulerprompt}
\begin{euleroutput}
  [4.71566,  0.38319]
\end{euleroutput}
\begin{eulercomment}
Kita memasukkan semuanya ke dalam satu plot untuk titik-titik data dan
garis-garis regresi, serta untuk bobot-bobot yang digunakan.
\end{eulercomment}
\begin{eulerprompt}
>figure(2,1);  ...
>figure(1); statplot(x,y,"b",xl="Regression"); ...
>  plot2d("evalpoly(x,p)",>add,color=blue,style="--"); ...
>  plot2d("evalpoly(x,pw)",5,10,>add,color=red,style="--"); ...
>figure(2); plot2d(w,1,10,>filled,style="/",fillcolor=red,xl=w); ...
>figure(0):
\end{eulerprompt}
\eulerimg{27}{Image Use/TugasAplikom8Adi25030630001-035.png}
\begin{eulercomment}
Sebagai contoh lainnya, kita membaca data survei mahasiswa—meliputi
usia mereka, usia orang tua mereka, serta jumlah saudara kandung—dari
sebuah berkas.

Tabel ini memuat "m" dan "f" pada kolom kedua. Kita menggunakan
variabel tok2 untuk menetapkan padanan terjemahan yang tepat alih-alih
membiarkan readtable() mengumpulkan terjemahannya sendiri.
\end{eulercomment}
\begin{eulerprompt}
>\{MS,hd\}:=readtable("table1.dat",tok2:=["m","f"]);  ...
>writetable(MS,labc=hd,tok2:=["m","f"]);
\end{eulerprompt}
\begin{euleroutput}
      Person       Sex       Age    Mother    Father  Siblings
           1         m        29        58        61         1
           2         f        26        53        54         2
           3         m        24        49        55         1
           4         f        25        56        63         3
           5         f        25        49        53         0
           6         f        23        55        55         2
           7         m        23        48        54         2
           8         m        27        56        58         1
           9         m        25        57        59         1
          10         m        24        50        54         1
          11         f        26        61        65         1
          12         m        24        50        52         1
          13         m        29        54        56         1
          14         m        28        48        51         2
          15         f        23        52        52         1
          16         m        24        45        57         1
          17         f        24        59        63         0
          18         f        23        52        55         1
          19         m        24        54        61         2
          20         f        23        54        55         1
\end{euleroutput}
\begin{eulercomment}
Bagaimana usia-usia tersebut saling bergantung satu sama lain?
Gambaran awal diperoleh melalui pairwise scatterplot.
\end{eulercomment}
\begin{eulerprompt}
>scatterplots(tablecol(MS,3:5),hd[3:5]):
\end{eulerprompt}
\eulerimg{27}{Image Use/TugasAplikom8Adi25030630001-036.png}
\begin{eulercomment}
Jelas terlihat bahwa usia ayah dan usia ibu saling bergantung satu
sama lain. Mari kita tentukan dan plot garis regresinya.
\end{eulercomment}
\begin{eulerprompt}
>cs:=MS[,4:5]'; ps:=polyfit(cs[1],cs[2],1)
\end{eulerprompt}
\begin{euleroutput}
  [17.3789,  0.740964]
\end{euleroutput}
\begin{eulercomment}
Model ini jelas keliru. Garis regresinya akan menjadi s=17+0.74t,
dengan t menyatakan usia ibu dan s menyatakan usia ayah. Selisih usia
mungkin sedikit bergantung pada usia, tetapi tidak sebesar itu.

Sebaliknya, kita menduga adanya fungsi seperti s=a+t. Maka a adalah
rerata dari s-t. Nilai tersebut merupakan selisih usia rata-rata
antara ayah dan ibu.
\end{eulercomment}
\begin{eulerprompt}
>da:=mean(cs[2]-cs[1])
\end{eulerprompt}
\begin{euleroutput}
  3.65
\end{euleroutput}
\begin{eulercomment}
Mari kita plot ini ke dalam satu plot sebar.
\end{eulercomment}
\begin{eulerprompt}
>plot2d(cs[1],cs[2],>points);  ...
>plot2d("evalpoly(x,ps)",color=red,style=".",>add);  ...
>plot2d("x+da",color=blue,>add):
\end{eulerprompt}
\eulerimg{27}{Image Use/TugasAplikom8Adi25030630001-037.png}
\begin{eulercomment}
Berikut adalah box plot dari kedua usia tersebut. Grafik ini hanya
menunjukkan bahwa kedua usia tersebut berbeda.
\end{eulercomment}
\begin{eulerprompt}
>boxplot(cs,["mothers","fathers"]):
\end{eulerprompt}
\eulerimg{27}{Image Use/TugasAplikom8Adi25030630001-038.png}
\begin{eulercomment}
Menariknya, selisih pada nilai-nilai median tidak sebesar selisih pada
nilai-nilai rerata.
\end{eulercomment}
\begin{eulerprompt}
>median(cs[2])-median(cs[1])
\end{eulerprompt}
\begin{euleroutput}
  1.5
\end{euleroutput}
\begin{eulercomment}
Koefisien korelasi menunjukkan adanya korelasi positif.
\end{eulercomment}
\begin{eulerprompt}
>correl(cs[1],cs[2])
\end{eulerprompt}
\begin{euleroutput}
  0.7588307236
\end{euleroutput}
\begin{eulercomment}
Korelasi peringkat merupakan ukuran kesamaan urutan pada kedua vektor.
Nilai ini juga cukup positif.
\end{eulercomment}
\begin{eulerprompt}
>rankcorrel(cs[1],cs[2])
\end{eulerprompt}
\begin{euleroutput}
  0.758925292358
\end{euleroutput}
\eulerheading{Creating new Functions}
\begin{eulercomment}
Tentu saja, bahasa EMT dapat digunakan untuk memprogram fungsi-fungsi
baru. Misalnya, kita mendefinisikan fungsi skewness.

\end{eulercomment}
\begin{eulerformula}
\[
\text{sk}(x) = \dfrac{\sqrt{n} \sum_i (x_i-m)^3}{\left(\sum_i (x_i-m)^2\right)^{3/2}}
\]
\end{eulerformula}
\begin{eulercomment}
dengan m menyatakan rerata dari x.
\end{eulercomment}
\begin{eulerprompt}
>function skew (x:vector) ...
\end{eulerprompt}
\begin{eulerudf}
  m=mean(x);
  return sqrt(cols(x))*sum((x-m)^3)/(sum((x-m)^2))^(3/2);
  endfunction
\end{eulerudf}
\begin{eulercomment}
Seperti yang terlihat, kita dapat dengan mudah memanfaatkan bahasa
matriks untuk memperoleh implementasi yang sangat singkat dan efisien.
Mari kita coba fungsi ini.
\end{eulercomment}
\begin{eulerprompt}
>data=normal(20); skew(normal(10))
\end{eulerprompt}
\begin{euleroutput}
  0.521806329961
\end{euleroutput}
\begin{eulercomment}
Berikut adalah fungsi lainnya, yang disebut koefisien skewness
Pearson.
\end{eulercomment}
\begin{eulerprompt}
>function skew1 (x) := 3*(mean(x)-median(x))/dev(x)
>skew1(data)
\end{eulerprompt}
\begin{euleroutput}
  0.573102925949
\end{euleroutput}
\eulerheading{Monte Carlo Simulation}
\begin{eulercomment}
Euler dapat digunakan untuk menyimulasikan kejadian-kejadian acak.
Kita telah melihat beberapa contoh sederhana di atas. Berikut adalah
contoh lainnya, yang menyimulasikan 1.000 kali lemparan 3 dadu dan
menanyakan distribusi dari jumlahnya.
\end{eulercomment}
\begin{eulerprompt}
>ds:=sum(intrandom(1000,3,6))';  fs=getmultiplicities(3:18,ds)
\end{eulerprompt}
\begin{euleroutput}
  [5,  12,  25,  49,  73,  94,  128,  141,  124,  116,  105,  62,  37,
  19,  7,  3]
\end{euleroutput}
\begin{eulercomment}
Kita sekarang dapat memplot ini.
\end{eulercomment}
\begin{eulerprompt}
>columnsplot(fs,lab=3:18):
\end{eulerprompt}
\eulerimg{27}{Image Use/TugasAplikom8Adi25030630001-040.png}
\begin{eulercomment}
Menentukan distribusi yang diharapkan tidaklah begitu mudah. Untuk
keperluan ini, kita menggunakan rekursi tingkat lanjut.

Fungsi berikut menghitung banyaknya cara bilangan k dapat dinyatakan
sebagai jumlah dari n bilangan dalam rentang 1 hingga m. Fungsi ini
bekerja secara rekursif dengan cara yang cukup jelas.
\end{eulercomment}
\begin{eulerprompt}
>function map countways (k; n, m) ...
\end{eulerprompt}
\begin{eulerudf}
    if n==1 then return k>=1 && k<=m
    else
      sum=0; 
      loop 1 to m; sum=sum+countways(k-#,n-1,m); end;
      return sum;
    end;
  endfunction
\end{eulerudf}
\begin{eulercomment}
Berikut adalah hasilnya untuk tiga kali lemparan dadu.
\end{eulercomment}
\begin{eulerprompt}
>countways(5:25,5,5)
\end{eulerprompt}
\begin{euleroutput}
  [1,  5,  15,  35,  70,  121,  185,  255,  320,  365,  381,  365,  320,
  255,  185,  121,  70,  35,  15,  5,  1]
\end{euleroutput}
\begin{eulerprompt}
>cw=countways(3:18,3,6)
\end{eulerprompt}
\begin{euleroutput}
  [1,  3,  6,  10,  15,  21,  25,  27,  27,  25,  21,  15,  10,  6,  3,
  1]
\end{euleroutput}
\begin{eulercomment}
Kita menambahkan nilai-nilai yang diharapkan ke dalam plot.
\end{eulercomment}
\begin{eulerprompt}
>plot2d(cw/6^3*1000,>add); plot2d(cw/6^3*1000,>points,>add):
\end{eulerprompt}
\eulerimg{27}{Image Use/TugasAplikom8Adi25030630001-041.png}
\begin{eulercomment}
Untuk simulasi lainnya, simpangan dari nilai rerata sebanyak \textdollar{}n\textdollar{}
peubah acak berdistribusi normal baku N(0,1) adalah 1/\textbackslash{}sqrt\{n\}.
\end{eulercomment}
\begin{eulerprompt}
>longformat; 1/sqrt(10)
\end{eulerprompt}
\begin{euleroutput}
  0.316227766017
\end{euleroutput}
\begin{eulercomment}
Mari kita periksa hal ini dengan sebuah simulasi. Kita membangkitkan
10.000 kali sebanyak 10 vektor acak.
\end{eulercomment}
\begin{eulerprompt}
>M=normal(10000,10); dev(mean(M)')
\end{eulerprompt}
\begin{euleroutput}
  0.318861419326
\end{euleroutput}
\begin{eulerprompt}
>plot2d(mean(M)',>distribution):
\end{eulerprompt}
\eulerimg{27}{Image Use/TugasAplikom8Adi25030630001-042.png}
\begin{eulercomment}
Median dari 10 bilangan acak berdistribusi normal baku memiliki
simpangan yang lebih besar.
\end{eulercomment}
\begin{eulerprompt}
>dev(median(M)')
\end{eulerprompt}
\begin{euleroutput}
  0.376651504162
\end{euleroutput}
\begin{eulercomment}
Karena kita dapat dengan mudah membangkitkan random walk, kita dapat
menyimulasikan proses Wiener. Kita mengambil 1.000 langkah dari 1.000
proses. Kita kemudian memplot simpangan baku dan rerata dari langkah
ke-\textdollar{}n\textdollar{} proses-proses tersebut bersama dengan nilai-nilai yang
diharapkan berwarna merah.
\end{eulercomment}
\begin{eulerprompt}
>n=1000; m=1000; M=cumsum(normal(n,m)/sqrt(m)); ...
>t=(1:n)/n; figure(2,1); ...
>figure(1); plot2d(t,mean(M')'); plot2d(t,0,color=red,>add); ...
>figure(2); plot2d(t,dev(M')'); plot2d(t,sqrt(t),color=red,>add); ...
>figure(0):
\end{eulerprompt}
\eulerimg{27}{Image Use/TugasAplikom8Adi25030630001-043.png}
\eulerheading{Tests}
\begin{eulercomment}
Uji statistik merupakan perangkat yang penting dalam statistika. Di
dalam Euler, banyak uji yang telah diimplementasikan. Seluruh uji
tersebut mengembalikan besar galat yang kita terima apabila kita
menolak hipotesis nol.

Sebagai contoh, kita menguji hasil lemparan dadu terhadap distribusi
seragam. Pada 600 kali lemparan, kita memperoleh nilai-nilai berikut,
yang kemudian kita masukkan ke dalam uji khi-kuadrat.
\end{eulercomment}
\begin{eulerprompt}
>chitest([90,103,114,101,103,89],dup(100,6)')
\end{eulerprompt}
\begin{euleroutput}
  0.498830517952
\end{euleroutput}
\begin{eulercomment}
Uji khi-kuadrat juga memiliki mode yang menggunakan simulasi Monte
Carlo untuk menguji statistiknya. Hasilnya seharusnya hampir sama.
Parameter \textgreater{}p menginterpretasikan vektor y sebagai vektor peluang.
\end{eulercomment}
\begin{eulerprompt}
>chitest([90,103,114,101,103,89],dup(1/6,6)',>p,>montecarlo)
\end{eulerprompt}
\begin{euleroutput}
  0.517
\end{euleroutput}
\begin{eulercomment}
Galat ini terlalu besar. Oleh karena itu, kita tidak dapat menolak
distribusi seragam. Hal ini tidak membuktikan bahwa dadu kita adil
(fair), tetapi kita tidak dapat menolak hipotesis kita.

Berikutnya, kita membangkitkan 1.000 kali lemparan dadu menggunakan
pembangkit bilangan acak, lalu melakukan uji yang sama.
\end{eulercomment}
\begin{eulerprompt}
>n=1000; t=random([1,n*6]); chitest(count(t*6,6),dup(n,6)')
\end{eulerprompt}
\begin{euleroutput}
  0.608613608204
\end{euleroutput}
\begin{eulercomment}
Mari kita uji untuk nilai rerata 100 dengan uji-t.
\end{eulercomment}
\begin{eulerprompt}
>s=200+normal([1,100])*10; ...
>ttest(mean(s),dev(s),100,200)
\end{eulerprompt}
\begin{euleroutput}
  0.39980320739
\end{euleroutput}
\begin{eulercomment}
Fungsi ttest() memerlukan nilai rerata, simpangan, banyaknya data,
serta nilai rerata yang akan diuji.

Sekarang mari kita periksa dua pengukuran untuk rerata yang sama. Kita
menolak hipotesis bahwa keduanya memiliki rerata yang sama jika
hasilnya kurang dari 0,05.
\end{eulercomment}
\begin{eulerprompt}
>tcomparedata(normal(1,10),normal(1,10))
\end{eulerprompt}
\begin{euleroutput}
  0.275550209428
\end{euleroutput}
\begin{eulercomment}
Jika kita menambahkan bias pada salah satu distribusi, kita memperoleh
lebih banyak penolakan. Ulangi simulasi ini beberapa kali untuk
melihat efeknya.
\end{eulercomment}
\begin{eulerprompt}
>tcomparedata(normal(1,10),normal(1,10)+2)
\end{eulerprompt}
\begin{euleroutput}
  3.76360968102e-06
\end{euleroutput}
\begin{eulercomment}
Pada contoh berikutnya, kita membangkitkan 20 lemparan dadu acak
sebanyak 100 kali dan menghitung banyaknya angka satu di dalamnya.
Rata-rata harus terdapat 20/6 = 3,3 angka satu.
\end{eulercomment}
\begin{eulerprompt}
>R=random(100,20); R=sum(R*6<=1)'; mean(R)
\end{eulerprompt}
\begin{euleroutput}
  3.14
\end{euleroutput}
\begin{eulercomment}
Kita sekarang membandingkan banyaknya angka satu dengan distribusi
binomial. Pertama, kita memplot distribusi dari banyaknya angka satu.
\end{eulercomment}
\begin{eulerprompt}
>plot2d(R,distribution=max(R)+1,even=1,style="\(\backslash\)/"):
\end{eulerprompt}
\eulerimg{27}{Image Use/TugasAplikom8Adi25030630001-044.png}
\begin{eulerprompt}
>t=count(R,21);
\end{eulerprompt}
\begin{eulercomment}
Kemudian kita menghitung nilai-nilai yang diharapkan.
\end{eulercomment}
\begin{eulerprompt}
>n=0:20; b=bin(20,n)*(1/6)^n*(5/6)^(20-n)*100;
\end{eulerprompt}
\begin{eulercomment}
Kita harus menggabungkan beberapa nilai untuk memperoleh
kategori-kategori yang cukup besar.
\end{eulercomment}
\begin{eulerprompt}
>t1=sum(t[1:2])|t[3:7]|sum(t[8:21]); ...
>b1=sum(b[1:2])|b[3:7]|sum(b[8:21]);
\end{eulerprompt}
\begin{eulercomment}
Uji khi-kuadrat menolak hipotesis bahwa distribusi kita merupakan
distribusi binomial jika hasilnya kurang dari 0,05.
\end{eulercomment}
\begin{eulerprompt}
>chitest(t1,b1)
\end{eulerprompt}
\begin{euleroutput}
  0.405510807841
\end{euleroutput}
\begin{eulercomment}
Contoh berikut memuat hasil dari dua kelompok orang (misalnya
laki-laki dan perempuan) yang memilih salah satu dari enam partai.
\end{eulercomment}
\begin{eulerprompt}
>A=[23,37,43,52,64,74;27,39,41,49,63,76];  ...
>  writetable(A,wc=6,labr=["m","f"],labc=1:6)
\end{eulerprompt}
\begin{euleroutput}
             1     2     3     4     5     6
       m    23    37    43    52    64    74
       f    27    39    41    49    63    76
\end{euleroutput}
\begin{eulercomment}
Kita ingin menguji independensi pilihan suara terhadap jenis kelamin.
Uji tabel \textdollar{}\textbackslash{}chi\textasciicircum{}2\textdollar{} melakukan pengujian ini. Hasilnya terlalu besar
untuk menolak independensi. Oleh karena itu, berdasarkan data ini,
kita tidak dapat menyatakan apakah pilihan suara bergantung pada jenis
kelamin.
\end{eulercomment}
\begin{eulerprompt}
>tabletest(A)
\end{eulerprompt}
\begin{euleroutput}
  0.990701632326
\end{euleroutput}
\begin{eulercomment}
Tabel berikut merupakan tabel ekspektasi apabila kita mengasumsikan
frekuensi teramati dari pilihan suara tersebut.
\end{eulercomment}
\begin{eulerprompt}
>writetable(expectedtable(A),wc=6,dc=1,labr=["m","f"],labc=1:6)
\end{eulerprompt}
\begin{euleroutput}
             1     2     3     4     5     6
       m  24.9  37.9  41.9  50.3  63.3  74.7
       f  25.1  38.1  42.1  50.7  63.7  75.3
\end{euleroutput}
\begin{eulercomment}
Kita dapat menghitung koefisien kontingensi terkoreksi. Karena
nilainya sangat dekat dengan 0, kita menyimpulkan bahwa pilihan suara
tidak bergantung pada jenis kelamin.
\end{eulercomment}
\begin{eulerprompt}
>contingency(A)
\end{eulerprompt}
\begin{euleroutput}
  0.0427225484717
\end{euleroutput}
\begin{eulercomment}
\begin{eulercomment}
\eulerheading{Some More Tests}
\begin{eulercomment}
Berikutnya, kita menggunakan analisis varians (uji-F) untuk menguji
tiga sampel data berdistribusi normal agar memiliki nilai rerata yang
sama. Metode ini disebut ANOVA (analysis of variance). Di Euler,
fungsi varanalysis() digunakan.
\end{eulercomment}
\begin{eulerprompt}
>x1=[109,111,98,119,91,118,109,99,115,109,94]; mean(x1),
\end{eulerprompt}
\begin{euleroutput}
  106.545454545
\end{euleroutput}
\begin{eulerprompt}
>x2=[120,124,115,139,114,110,113,120,117]; mean(x2),
\end{eulerprompt}
\begin{euleroutput}
  119.111111111
\end{euleroutput}
\begin{eulerprompt}
>x3=[120,112,115,110,105,134,105,130,121,111]; mean(x3)
\end{eulerprompt}
\begin{euleroutput}
  116.3
\end{euleroutput}
\begin{eulerprompt}
>varanalysis(x1,x2,x3)
\end{eulerprompt}
\begin{euleroutput}
  0.0138048221371
\end{euleroutput}
\begin{eulercomment}
Ini berarti kita menolak hipotesis bahwa nilai reratanya sama. Kita
melakukan hal ini dengan peluang galat sebesar 1,3\%.

Terdapat pula uji median, yang menolak sampel-sampel data dengan
sebaran rerata yang berbeda dengan menguji median dari sampel gabungan
tersebut.
\end{eulercomment}
\begin{eulerprompt}
>a=[56,66,68,49,61,53,45,58,54];
>b=[72,81,51,73,69,78,59,67,65,71,68,71];
>mediantest(a,b)
\end{eulerprompt}
\begin{euleroutput}
  0.0241724220052
\end{euleroutput}
\begin{eulercomment}
Uji lain mengenai kesamaan adalah uji peringkat. Uji ini jauh lebih
tajam dibandingkan uji median.
\end{eulercomment}
\begin{eulerprompt}
>ranktest(a,b)
\end{eulerprompt}
\begin{euleroutput}
  0.00199969612469
\end{euleroutput}
\begin{eulercomment}
Pada contoh berikut, kedua distribusi memiliki rerata yang sama.
\end{eulercomment}
\begin{eulerprompt}
>ranktest(random(1,100),random(1,50)*3-1)
\end{eulerprompt}
\begin{euleroutput}
  0.485688783753
\end{euleroutput}
\begin{eulercomment}
Sekarang mari kita coba mensimulasikan dua perlakuan a dan b yang
diterapkan pada orang-orang yang berbeda.
\end{eulercomment}
\begin{eulerprompt}
>a=[8.0,7.4,5.9,9.4,8.6,8.2,7.6,8.1,6.2,8.9];
>b=[6.8,7.1,6.8,8.3,7.9,7.2,7.4,6.8,6.8,8.1];
\end{eulerprompt}
\begin{eulercomment}
Uji tanda menentukan apakah a lebih baik daripada b.
\end{eulercomment}
\begin{eulerprompt}
>signtest(a,b)
\end{eulerprompt}
\begin{euleroutput}
  0.0546875
\end{euleroutput}
\begin{eulercomment}
Galat ini terlalu besar. Kita tidak dapat menolak hipotesis bahwa a
sama baiknya dengan b.

Uji Wilcoxon lebih tajam daripada uji ini, tetapi bergantung pada
nilai kuantitatif dari selisih-selisih tersebut.
\end{eulercomment}
\begin{eulerprompt}
>wilcoxon(a,b)
\end{eulerprompt}
\begin{euleroutput}
  0.0296680599405
\end{euleroutput}
\begin{eulercomment}
Mari kita coba dua uji lagi dengan menggunakan deret yang
dibangkitkan.
\end{eulercomment}
\begin{eulerprompt}
>wilcoxon(normal(1,20),normal(1,20)-1)
\end{eulerprompt}
\begin{euleroutput}
  0.0284572962926
\end{euleroutput}
\begin{eulerprompt}
>wilcoxon(normal(1,20),normal(1,20))
\end{eulerprompt}
\begin{euleroutput}
  0.87642769873
\end{euleroutput}
\eulerheading{Random Numbers}
\begin{eulercomment}
Berikut ini adalah pengujian untuk generator angka acak. Euler
menggunakan generator yang sangat baik, jadi kita tidak perlu
mengkhawatirkan adanya masalah.

Pertama, kita menghasilkan sepuluh juta angka acak dalam interval
[0,1].
\end{eulercomment}
\begin{eulerprompt}
>n:=10000000; r:=random(1,n);
\end{eulerprompt}
\begin{eulercomment}
Selanjutnya, kita menghitung jarak antara dua angka yang kurang dari
0,05.
\end{eulercomment}
\begin{eulerprompt}
>a:=0.05; d:=differences(nonzeros(r<a));
\end{eulerprompt}
\begin{eulercomment}
Terakhir, kita memplot berapa kali setiap jarak muncul, dan
membandingkannya dengan nilai yang diharapkan.
\end{eulercomment}
\begin{eulerprompt}
>m=getmultiplicities(1:100,d); plot2d(m); ...
>  plot2d("n*(1-a)^(x-1)*a^2",color=red,>add):
\end{eulerprompt}
\eulerimg{27}{Image Use/TugasAplikom8Adi25030630001-045.png}
\begin{eulercomment}
Bersihkan data.
\end{eulercomment}
\begin{eulerprompt}
>remvalue n;
\end{eulerprompt}
\begin{eulercomment}
\begin{eulercomment}
\eulerheading{Introduction for Users of the R Project}
\begin{eulercomment}
Jelas bahwa EMT tidak bersaing dengan R sebagai paket statistik.
Namun, ada banyak prosedur dan fungsi statistik yang juga tersedia di
EMT. Jadi, EMT dapat memenuhi kebutuhan dasar. Lagipula, EMT
dilengkapi dengan paket numerik dan sistem aljabar komputer.

Catatan buku kerja (notebook) ini ditujukan bagi Anda yang sudah
familiar dengan R, tetapi perlu mengetahui perbedaan sintaksis antara
EMT dan R. Kami berupaya memberikan gambaran tentang hal-hal yang
jelas maupun yang kurang jelas yang perlu Anda ketahui.

Selain itu, kita akan melihat cara bertukar data antara kedua sistem
tersebut
\end{eulercomment}
\begin{eulercomment}
Perlu dicatat bahwa ini masih dalam proses pengerjaan.
\end{eulercomment}
\eulerheading{Basic Syntax}
\begin{eulercomment}
Hal pertama yang Anda pelajari di R adalah membuat vektor. Di EMT,
perbedaan utamanya adalah operator : dapat menerima ukuran langkah
(step size). Selain itu, operator ini memiliki prioritas pengikatan
(binding power) yang rendah.
\end{eulercomment}
\begin{eulerprompt}
>n=10; 0:n/20:n-1
\end{eulerprompt}
\begin{euleroutput}
  [0,  0.5,  1,  1.5,  2,  2.5,  3,  3.5,  4,  4.5,  5,  5.5,  6,  6.5,
  7,  7.5,  8,  8.5,  9]
\end{euleroutput}
\begin{eulercomment}
Fungsi c() tidak ada. Vektor dapat digunakan untuk menggabungkan
berbagai elemen.

Contoh berikut, seperti banyak contoh lainnya, berasal dari
"Introduction to R" yang disertakan bersama proyek R. Jika Anda
membaca PDF ini, Anda akan menemukan bahwa saya mengikuti alurnya
dalam tutorial ini.
\end{eulercomment}
\begin{eulerprompt}
>x=[10.4, 5.6, 3.1, 6.4, 21.7]; [x,0,x]
\end{eulerprompt}
\begin{euleroutput}
  [10.4,  5.6,  3.1,  6.4,  21.7,  0,  10.4,  5.6,  3.1,  6.4,  21.7]
\end{euleroutput}
\begin{eulercomment}
Operator titik dua dengan ukuran langkah di EMT digantikan oleh fungsi
seq() di R. Kita dapat menulis fungsi ini di EMT.
\end{eulercomment}
\begin{eulerprompt}
>function seq(a,b,c) := a:b:c; ...
>seq(0,-0.1,-1)
\end{eulerprompt}
\begin{euleroutput}
  [0,  -0.1,  -0.2,  -0.3,  -0.4,  -0.5,  -0.6,  -0.7,  -0.8,  -0.9,  -1]
\end{euleroutput}
\begin{eulercomment}
Fungsi rep() dari R tidak tersedia di EMT. Untuk input vektor, fungsi
ini dapat ditulis sebagai berikut.
\end{eulercomment}
\begin{eulerprompt}
>function rep(x:vector,n:index) := flatten(dup(x,n)); ...
>rep(x,2)
\end{eulerprompt}
\begin{euleroutput}
  [10.4,  5.6,  3.1,  6.4,  21.7,  10.4,  5.6,  3.1,  6.4,  21.7]
\end{euleroutput}
\begin{eulercomment}
Perlu dicatat bahwa = atau := digunakan untuk penugasan (assignment).
Operator -\textgreater{} digunakan untuk satuan (units) di EMT.
\end{eulercomment}
\begin{eulerprompt}
>125km -> " miles"
\end{eulerprompt}
\begin{euleroutput}
  77.6713990297 miles
\end{euleroutput}
\begin{eulercomment}
Operator "\textless{}-" untuk penugasan memang membingungkan, dan bukan ide yang
bagus dari R. Perintah berikut akan membandingkan a dan -4 di EMT.
\end{eulercomment}
\begin{eulerprompt}
>a=2; a<-4
\end{eulerprompt}
\begin{euleroutput}
  0
\end{euleroutput}
\begin{eulercomment}
Di R, a\textless{}-4\textless{}3 dapat berfungsi, tetapi a\textless{}-4\textless{}-3 tidak. Saya juga
mengalami ambiguitas serupa di EMT, namun berusaha menghilangkannya
secara bertahap.

EMT dan R memiliki vektor bertipe boolean. Namun di EMT, angka 0 dan 1
digunakan untuk merepresentasikan false dan true. Meskipun demikian,
di R nilai true dan false tetap dapat digunakan dalam aritmetika biasa
seperti halnya di EMT.
\end{eulercomment}
\begin{eulerprompt}
>x<5, %*x
\end{eulerprompt}
\begin{euleroutput}
  [0,  0,  1,  0,  0]
  [0,  0,  3.1,  0,  0]
\end{euleroutput}
\begin{eulercomment}
EMT memunculkan kesalahan (error) atau menghasilkan NAN bergantung
pada penanda (flag) 'errors'.
\end{eulercomment}
\begin{eulerprompt}
>errors off; 0/0, isNAN(sqrt(-1)), errors on;
\end{eulerprompt}
\begin{euleroutput}
  NAN
  1
\end{euleroutput}
\begin{eulercomment}
String di R dan EMT adalah sama. Keduanya menggunakan locale
(pengaturan lokal) saat ini, bukan Unicode.

Di R, terdapat paket-paket untuk Unicode. Di EMT, sebuah string dapat
berupa string Unicode. String Unicode dapat diterjemahkan ke pengodean
lokal (local encoding) dan sebaliknya. Selain itu, u"..." dapat memuat
entitas HTML.
\end{eulercomment}
\begin{eulerprompt}
>u"&#169; Ren&eacut; Grothmann"
\end{eulerprompt}
\begin{euleroutput}
  © René Grothmann
\end{euleroutput}
\begin{eulercomment}
Karakter berikut mungkin dapat atau tidak dapat ditampilkan dengan
benar di sistem Anda sebagai huruf A dengan titik dan garis di
atasnya. Hal ini bergantung pada font yang Anda gunakan.
\end{eulercomment}
\begin{eulerprompt}
>chartoutf([480])
\end{eulerprompt}
\begin{euleroutput}
  Ç 
\end{euleroutput}
\begin{eulercomment}
Penggabungan string (string concatenation) dilakukan dengan + atau \textbar{}.
Penggabungan ini dapat menyertakan angka, yang akan dicetak dalam
format saat ini.
\end{eulercomment}
\begin{eulerprompt}
>"pi = "+pi
\end{eulerprompt}
\begin{euleroutput}
  pi = 3.14159265359
\end{euleroutput}
\eulerheading{Indexing}
\begin{eulercomment}
Sebagian besar waktu, hal ini akan berfungsi sama seperti di R.

Namun, EMT akan menafsirkan indeks negatif dari bagian belakang
vektor, sedangkan R menafsirkan x[-n] (tertulis x[n] pada teks asli)
sebagai x tanpa elemen ke-n.
\end{eulercomment}
\begin{eulerprompt}
>x, x[1:3], x[-2]
\end{eulerprompt}
\begin{euleroutput}
  [10.4,  5.6,  3.1,  6.4,  21.7]
  [10.4,  5.6,  3.1]
  6.4
\end{euleroutput}
\begin{eulercomment}
Perilaku R tersebut dapat dicapai di EMT dengan menggunakan drop().
\end{eulercomment}
\begin{eulerprompt}
>drop(x,2)
\end{eulerprompt}
\begin{euleroutput}
  [10.4,  3.1,  6.4,  21.7]
\end{euleroutput}
\begin{eulercomment}
Vektor logika tidak diperlakukan secara berbeda sebagai indeks di EMT,
berbeda dengan di R. Anda perlu mengekstrak elemen bukan nol
(non-zero) terlebih dahulu di EMT.
\end{eulercomment}
\begin{eulerprompt}
>x, x>5, x[nonzeros(x>5)]
\end{eulerprompt}
\begin{euleroutput}
  [10.4,  5.6,  3.1,  6.4,  21.7]
  [1,  1,  0,  1,  1]
  [10.4,  5.6,  6.4,  21.7]
\end{euleroutput}
\begin{eulercomment}
Sama seperti di R, vektor indeks dapat berisi pengulangan.
\end{eulercomment}
\begin{eulerprompt}
>x[[1,2,2,1]]
\end{eulerprompt}
\begin{euleroutput}
  [10.4,  5.6,  5.6,  10.4]
\end{euleroutput}
\begin{eulercomment}
Namun, penamaan untuk indeks tidak dimungkinkan di EMT. Untuk sebuah
paket statistik, hal ini sering kali diperlukan guna mempermudah akses
ke elemen-elemen vektor.

Untuk meniru perilaku ini, kita dapat mendefinisikan sebuah fungsi
seperti berikut.
\end{eulercomment}
\begin{eulerprompt}
>function pilih (v,i,s) := v[indexof(s,i)]; ...
> s=["first","second","third","fourth"]; pilih(x,["first","third"],s)
\end{eulerprompt}
\begin{euleroutput}
  [10.4,  3.1]
\end{euleroutput}
\eulerheading{Data Types}
\begin{eulercomment}
EMT memiliki tipe data yang lebih tetap (fixed) dibandingkan R.
Jelasnya, di R terdapat vektor yang dapat tumbuh (growing vectors).
Anda dapat mengatur vektor numerik kosong v dan menetapkan nilai ke
elemen v[17]. Hal ini tidak dimungkinkan di EMT.

Contoh berikut ini agak kurang efisien.
\end{eulercomment}
\begin{eulerprompt}
>v=[]; for i=1 to 10000; v=v|i; end;
\end{eulerprompt}
\begin{eulercomment}
EMT kini akan menyusun sebuah vektor dengan v dan i yang ditambahkan
pada stack lalu menyalin kembali vektor tersebut ke variabel global v.

Cara yang lebih efisien adalah dengan mendefinisikan vektor terlebih
dahulu (pre-defines).
\end{eulercomment}
\begin{eulerprompt}
>v=zeros(10000); for i=1 to 10000; v[i]=i; end;
\end{eulerprompt}
\begin{eulercomment}
Untuk mengubah tipe data di EMT, Anda dapat menggunakan fungsi seperti
complex().
\end{eulercomment}
\begin{eulerprompt}
>complex(1:4)
\end{eulerprompt}
\begin{euleroutput}
  [ 1+0i ,  2+0i ,  3+0i ,  4+0i  ]
\end{euleroutput}
\begin{eulercomment}
Konversi ke string hanya dimungkinkan untuk tipe data elementer.
Format saat ini digunakan untuk penggabungan string sederhana. Namun,
terdapat fungsi seperti print() atau frac().

Untuk vektor, Anda dapat dengan mudah menulis fungsi Anda sendiri.
\end{eulercomment}
\begin{eulerprompt}
>function tostr (v) ...
\end{eulerprompt}
\begin{eulerudf}
  s="[";
  loop 1 to length(v);
     s=s+print(v[#],2,0);
     if #<length(v) then s=s+","; endif;
  end;
  return s+"]";
  endfunction
\end{eulerudf}
\begin{eulerprompt}
>tostr(linspace(0,1,10))
\end{eulerprompt}
\begin{euleroutput}
  [0.00,0.10,0.20,0.30,0.40,0.50,0.60,0.70,0.80,0.90,1.00]
\end{euleroutput}
\begin{eulercomment}
Untuk komunikasi dengan Maxima, terdapat fungsi convertmxm(), yang
juga dapat digunakan untuk memformat vektor untuk keluaran (output).
\end{eulercomment}
\begin{eulerprompt}
>convertmxm(1:10)
\end{eulerprompt}
\begin{euleroutput}
  [1,2,3,4,5,6,7,8,9,10]
\end{euleroutput}
\begin{eulercomment}
Untuk LaTeX, perintah tex dapat digunakan untuk memperoleh
kode/perintah LaTeX.
\end{eulercomment}
\begin{eulerprompt}
>tex(&[1,2,3])
\end{eulerprompt}
\begin{euleroutput}
  \(\backslash\)left[ 1 , 2 , 3 \(\backslash\)right] 
\end{euleroutput}
\eulerheading{Factors and Tables}
\begin{eulercomment}
Dalam pengantar R terdapat contoh dengan yang disebut faktor
(factors).

Berikut ini adalah daftar wilayah dari 30 negara bagian.
\end{eulercomment}
\begin{eulerprompt}
>austates = ["tas", "sa", "qld", "nsw", "nsw", "nt", "wa", "wa", ...
>"qld", "vic", "nsw", "vic", "qld", "qld", "sa", "tas", ...
>"sa", "nt", "wa", "vic", "qld", "nsw", "nsw", "wa", ...
>"sa", "act", "nsw", "vic", "vic", "act"];
\end{eulerprompt}
\begin{eulercomment}
Asumsikan kita memiliki pendapatan yang bersesuaian di setiap negara
bagian.
\end{eulercomment}
\begin{eulerprompt}
>incomes = [60, 49, 40, 61, 64, 60, 59, 54, 62, 69, 70, 42, 56, ...
>61, 61, 61, 58, 51, 48, 65, 49, 49, 41, 48, 52, 46, ...
>59, 46, 58, 43];
\end{eulerprompt}
\begin{eulercomment}
Kini, kita ingin menghitung rata-rata pendapatan di masing-masing
wilayah. Sebagai program statistik, R memiliki factor() dan tapply()
(tertulis tappy() pada teks asli) untuk keperluan ini.

EMT dapat melakukan hal ini dengan mencari indeks wilayah di dalam
daftar wilayah yang unik (unique list of territories).
\end{eulercomment}
\begin{eulerprompt}
>auterr=sort(unique(austates)); f=indexofsorted(auterr,austates)
\end{eulerprompt}
\begin{euleroutput}
  [6,  5,  4,  2,  2,  3,  8,  8,  4,  7,  2,  7,  4,  4,  5,  6,  5,  3,
  8,  7,  4,  2,  2,  8,  5,  1,  2,  7,  7,  1]
\end{euleroutput}
\begin{eulercomment}
Pada titik tersebut, kita dapat menulis fungsi perulangan (loop) kita
sendiri untuk melakukan operasi hanya pada satu faktor.

Atau kita dapat meniru fungsi tapply() dengan cara berikut.
\end{eulercomment}
\begin{eulerprompt}
>function map tappl (i; f$:call, cat, x) ...
\end{eulerprompt}
\begin{eulerudf}
  u=sort(unique(cat));
  f=indexof(u,cat);
  return f$(x[nonzeros(f==indexof(u,i))]);
  endfunction
\end{eulerudf}
\begin{eulercomment}
Hal ini agak kurang efisien karena menghitung wilayah yang unik untuk
setiap i, namun tetap berfungsi dengan baik.
\end{eulercomment}
\begin{eulerprompt}
>tappl(auterr,"mean",austates,incomes)
\end{eulerprompt}
\begin{euleroutput}
  [44.5,  57.3333,  55.5,  53.6,  55,  60.5,  56,  52.25]
\end{euleroutput}
\begin{eulercomment}
Perhatikan bahwa cara ini berlaku untuk setiap vektor wilayah.
\end{eulercomment}
\begin{eulerprompt}
>tappl(["act","nsw"],"mean",austates,incomes)
\end{eulerprompt}
\begin{euleroutput}
  [44.5,  57.3333]
\end{euleroutput}
\begin{eulercomment}
Kini, paket statistik EMT mendefinisikan tabel (tables) sama seperti
di R. Fungsi readtable() dan writetable() dapat digunakan untuk input
dan output.

Dengan demikian, kita dapat menampilkan rata-rata pendapatan negara
bagian di tiap wilayah dengan cara yang mudah dibaca (friendly).
\end{eulercomment}
\begin{eulerprompt}
>writetable(tappl(auterr,"mean",austates,incomes),labc=auterr,wc=7)
\end{eulerprompt}
\begin{euleroutput}
      act    nsw     nt    qld     sa    tas    vic     wa
     44.5  57.33   55.5   53.6     55   60.5     56  52.25
\end{euleroutput}
\begin{eulercomment}
Kita juga dapat mencoba meniru perilaku R secara lengkap.

Faktor-faktor tersebut sebaiknya disimpan secara jelas dalam sebuah
koleksi bersama tipe dan kategorinya (negara bagian dan wilayah pada
contoh kita). Untuk EMT, kita menambahkan indeks yang telah dihitung
sebelumnya (pre-computed indices).
\end{eulercomment}
\begin{eulerprompt}
>function makef (t) ...
\end{eulerprompt}
\begin{eulerudf}
  ## Factor data
  ## Returns a collection with data t, unique data, indices.
  ## See: tapply
  u=sort(unique(t));
  return \{\{t,u,indexofsorted(u,t)\}\};
  endfunction
\end{eulerudf}
\begin{eulerprompt}
>statef=makef(austates);
\end{eulerprompt}
\begin{eulercomment}
Kini elemen ketiga dari koleksi tersebut akan berisi indeks-indeksnya.
\end{eulercomment}
\begin{eulerprompt}
>statef[3]
\end{eulerprompt}
\begin{euleroutput}
  [6,  5,  4,  2,  2,  3,  8,  8,  4,  7,  2,  7,  4,  4,  5,  6,  5,  3,
  8,  7,  4,  2,  2,  8,  5,  1,  2,  7,  7,  1]
\end{euleroutput}
\begin{eulercomment}
Kini kita dapat meniru tapply() dengan cara berikut. Fungsi ini akan
mengembalikan sebuah tabel sebagai koleksi dari data tabel dan judul
kolom (column headings).
\end{eulercomment}
\begin{eulerprompt}
>function tapply (t:vector,tf,f$:call) ...
\end{eulerprompt}
\begin{eulerudf}
  ## Makes a table of data and factors
  ## tf : output of makef()
  ## See: makef
  uf=tf[2]; f=tf[3]; x=zeros(length(uf));
  for i=1 to length(uf);
     ind=nonzeros(f==i);
     if length(ind)==0 then x[i]=NAN;
     else x[i]=f$(t[ind]);
     endif;
  end;
  return \{\{x,uf\}\};
  endfunction
\end{eulerudf}
\begin{eulercomment}
Kita tidak menambahkan banyak pemeriksaan tipe di sini. Satu-satunya
pencegahan menyangkut kategori (faktor) yang tidak memiliki data.
Namun, sebaiknya diperiksa kesesuaian panjang t serta kebenaran
koleksi tf.

Tabel ini dapat dicetak sebagai tabel menggunakan writetable().
\end{eulercomment}
\begin{eulerprompt}
>writetable(tapply(incomes,statef,"mean"),wc=7)
\end{eulerprompt}
\begin{euleroutput}
      act    nsw     nt    qld     sa    tas    vic     wa
     44.5  57.33   55.5   53.6     55   60.5     56  52.25
\end{euleroutput}
\eulerheading{Arrays}
\begin{eulercomment}
EMT hanya memiliki dua dimensi untuk larik (arrays). Tipe datanya
disebut matriks (matrix). Namun, akan mudah untuk menulis fungsi untuk
dimensi yang lebih tinggi atau pustaka C untuk keperluan ini.

R memiliki lebih dari dua dimensi. Di R, larik adalah vektor dengan
atribut/bidang dimensi (dimension field).

Di EMT, vektor adalah matriks dengan satu baris. Vektor dapat diubah
menjadi matriks menggunakan redim().
\end{eulercomment}
\begin{eulerprompt}
>shortformat; X=redim(1:20,4,5)
\end{eulerprompt}
\begin{euleroutput}
          1         2         3         4         5 
          6         7         8         9        10 
         11        12        13        14        15 
         16        17        18        19        20 
\end{euleroutput}
\begin{eulercomment}
Ekstraksi baris dan kolom, atau sub-matriks, sangat mirip dengan yang
ada di R.
\end{eulercomment}
\begin{eulerprompt}
>X[,2:3]
\end{eulerprompt}
\begin{euleroutput}
          2         3 
          7         8 
         12        13 
         17        18 
\end{euleroutput}
\begin{eulercomment}
Namun, di R dimungkinkan untuk menetapkan nilai pada daftar indeks
tertentu dari vektor tersebut. Hal yang sama di EMT hanya dapat
dilakukan menggunakan perulangan (loop).
\end{eulercomment}
\begin{eulerprompt}
>function setmatrixvalue (M, i, j, v) ...
\end{eulerprompt}
\begin{eulerudf}
  loop 1 to max(length(i),length(j),length(v))
     M[i\{#\},j\{#\}] = v\{#\};
  end;
  endfunction
\end{eulerudf}
\begin{eulercomment}
Kita mendemonstrasikan hal ini untuk menunjukkan bahwa matriks
dilewatkan secara referensi (passed by reference) di EMT. Jika Anda
tidak ingin mengubah matriks asli M, Anda perlu menyalinnya di dalam
fungsi tersebut.
\end{eulercomment}
\begin{eulerprompt}
>setmatrixvalue(X,1:3,3:-1:1,0); X,
\end{eulerprompt}
\begin{euleroutput}
          1         2         0         4         5 
          6         0         8         9        10 
          0        12        13        14        15 
         16        17        18        19        20 
\end{euleroutput}
\begin{eulercomment}
Produk luar (outer product) di EMT hanya dapat dilakukan antarvektor.
Hal ini terjadi secara otomatis karena bahasa matriksnya. Satu vektor
harus berupa vektor kolom dan vektor lainnya berupa vektor baris.
\end{eulercomment}
\begin{eulerprompt}
>(1:5)*(1:5)'
\end{eulerprompt}
\begin{euleroutput}
          1         2         3         4         5 
          2         4         6         8        10 
          3         6         9        12        15 
          4         8        12        16        20 
          5        10        15        20        25 
\end{euleroutput}
\begin{eulercomment}
Dalam PDF pengantar untuk R terdapat sebuah contoh yang menghitung
distribusi dari ab-cd untuk a, b, c, d yang dipilih secara acak dari 0
hingga n. Solusi di R adalah membentuk matriks berdimensi 4 dan
menjalankan table() di atasnya.

Tentu saja, hal ini dapat dicapai dengan perulangan (loop). Namun,
perulangan kurang efektif di EMT maupun di R. Di EMT, kita bisa
menulis perulangan tersebut dalam bahasa C dan itu akan menjadi solusi
tercepat.

Akan tetapi, kita ingin meniru perilaku R. Untuk itu, kita perlu
meratakan (flatten) hasil perkalian ab dan membuat matriks dari ab-cd.
\end{eulercomment}
\begin{eulerprompt}
>a=0:6; b=a'; p=flatten(a*b); q=flatten(p-p'); ...
>u=sort(unique(q)); f=getmultiplicities(u,q); ...
>statplot(u,f,"h"):
\end{eulerprompt}
\eulerimg{27}{Image Use/TugasAplikom8Adi25030630001-046.png}
\begin{eulercomment}
Selain multiplisitas eksak (exact multiplicities), EMT dapat
menghitung frekuensi di dalam vektor.
\end{eulercomment}
\begin{eulerprompt}
>getfrequencies(q,-50:10:50)
\end{eulerprompt}
\begin{euleroutput}
  [0,  23,  132,  316,  602,  801,  333,  141,  53,  0]
\end{euleroutput}
\begin{eulercomment}
Cara termudah untuk memplot ini sebagai sebuah distribusi adalah
sebagai berikut.
\end{eulercomment}
\begin{eulerprompt}
>plot2d(q,distribution=11):
\end{eulerprompt}
\eulerimg{27}{Image Use/TugasAplikom8Adi25030630001-047.png}
\begin{eulercomment}
Namun, dimungkinkan juga untuk menghitung jumlah (count) pada interval
yang dipilih terlebih dahulu (beforehand). Tentu saja, cara berikut
menggunakan getfrequencies() secara internal.

Karena fungsi histo() mengembalikan frekuensi, kita perlu
menskalakannya agar integral di bawah grafik batang (bar graph)
bernilai 1.
\end{eulercomment}
\begin{eulerprompt}
>\{x,y\}=histo(q,v=-55:10:55); y=y/sum(y)/differences(x); ...
>plot2d(x,y,>bar,style="/"):
\end{eulerprompt}
\eulerimg{27}{Image Use/TugasAplikom8Adi25030630001-048.png}
\eulerheading{Lists}
\begin{eulercomment}
EMT memiliki dua jenis daftar (lists). Satu berupa daftar global yang
bersifat mutable (dapat diubah), dan lainnya berupa tipe daftar yang
bersifat immutable (tidak dapat diubah). Kita tidak membahas daftar
global di sini.

Tipe daftar yang immutable disebut collection (koleksi) di EMT. Tipe
ini berperilaku seperti struktur (struct) di C, tetapi
elemen-elemennya hanya diberi nomor dan tidak diberi nama.
\end{eulercomment}
\begin{eulerprompt}
>L=\{\{"Fred","Flintstone",40,[1990,1992]\}\}
\end{eulerprompt}
\begin{euleroutput}
  Fred
  Flintstone
  40
  [1990,  1992]
\end{euleroutput}
\begin{eulercomment}
Saat ini elemen-elemen tersebut tidak memiliki nama, meskipun nama
dapat diatur untuk tujuan khusus. Elemen-elemen tersebut diakses
menggunakan nomor.
\end{eulercomment}
\begin{eulerprompt}
>(L[4])[2]
\end{eulerprompt}
\begin{euleroutput}
  1992
\end{euleroutput}
\begin{eulercomment}
\begin{eulercomment}
\eulerheading{File Input and Output (Reading and Writing Data)}
\begin{eulercomment}
Anda akan sering ingin mengimpor matriks data dari sumber lain ke EMT.
Tutorial ini menjelaskan berbagai cara untuk mewujudkannya. Fungsi
sederhananya adalah writematrix() dan readmatrix().

Mari kita demonstrasikan cara membaca dan menulis vektor bilangan riil
ke sebuah berkas.
\end{eulercomment}
\begin{eulerprompt}
>a=random(1,100); mean(a), dev(a),
\end{eulerprompt}
\begin{euleroutput}
  0.47466
  0.27327
\end{euleroutput}
\begin{eulercomment}
Untuk menulis data ke sebuah berkas, kita menggunakan fungsi
writematrix().

Karena pengantar ini kemungkinan besar berada di direktori tempat
pengguna tidak memiliki akses tulis (write access), kita menulis data
tersebut ke direktori home pengguna. Untuk notebook milik sendiri, hal
ini tidak diperlukan karena berkas data akan ditulis ke direktori yang
sama.
\end{eulercomment}
\begin{eulerprompt}
>filename="test.dat";
\end{eulerprompt}
\begin{eulercomment}
Kini kita menulis vektor kolom a' ke berkas tersebut. Hal ini
menghasilkan satu angka di setiap baris berkas.
\end{eulercomment}
\begin{eulerprompt}
>writematrix(a',filename);
\end{eulerprompt}
\begin{eulercomment}
Untuk membaca data, kita menggunakan readmatrix().
\end{eulercomment}
\begin{eulerprompt}
>a=readmatrix(filename)';
\end{eulerprompt}
\begin{eulercomment}
Lalu hapus berkas tersebut.
\end{eulercomment}
\begin{eulerprompt}
>fileremove(filename);
>mean(a), dev(a),
\end{eulerprompt}
\begin{euleroutput}
  0.47466
  0.27327
\end{euleroutput}
\begin{eulercomment}
Fungsi writematrix() atau writetable() dapat dikonfigurasi untuk
bahasa lain.

Misalnya, jika Anda menggunakan sistem bahasa Indonesia (tanda desimal
berupa koma), Excel Anda membutuhkan nilai dengan koma desimal yang
dipisahkan oleh titik koma (semicolons) di dalam berkas csv (comma
separated values, yang standarnya dipisahkan oleh koma). Berkas
"test.csv" berikut seharusnya muncul di folder Anda saat ini.
\end{eulercomment}
\begin{eulerprompt}
>filename="test.csv"; ...
>writematrix(random(5,3),file=filename,separator=",");
\end{eulerprompt}
\begin{eulercomment}
Kini Anda dapat membuka berkas ini secara langsung menggunakan Excel
versi bahasa Indonesia.
\end{eulercomment}
\begin{eulerprompt}
>fileremove(filename);
\end{eulerprompt}
\begin{eulercomment}
Terkadang kita memiliki string dengan token seperti berikut ini.
\end{eulercomment}
\begin{eulerprompt}
>s1:="f m m f m m m f f f m m f";  ...
>s2:="f f f m m f f";
\end{eulerprompt}
\begin{eulercomment}
Untuk melakukan tokenisasi (tokenize) pada hal ini, kita
mendefinisikan sebuah vektor token.
\end{eulercomment}
\begin{eulerprompt}
>tok:=["f","m"]
\end{eulerprompt}
\begin{euleroutput}
  f
  m
\end{euleroutput}
\begin{eulercomment}
Lalu kita dapat menghitung berapa kali setiap token muncul di dalam
string tersebut, dan memasukkan hasilnya ke dalam sebuah tabel.
\end{eulercomment}
\begin{eulerprompt}
>M:=getmultiplicities(tok,strtokens(s1))_ ...
>  getmultiplicities(tok,strtokens(s2));
\end{eulerprompt}
\begin{eulercomment}
Tuliskan tabel tersebut dengan header berupa token.
\end{eulercomment}
\begin{eulerprompt}
>writetable(M,labc=tok,labr=1:2,wc=8)
\end{eulerprompt}
\begin{euleroutput}
                 f       m
         1       6       7
         2       5       2
\end{euleroutput}
\begin{eulercomment}
Untuk statistika, EMT dapat membaca dan menulis tabel.
\end{eulercomment}
\begin{eulerprompt}
>file="test.dat"; open(file,"w"); ...
>writeln("A,B,C"); writematrix(random(3,3)); ...
>close();
\end{eulerprompt}
\begin{eulercomment}
Berkas tersebut terlihat seperti ini.
\end{eulercomment}
\begin{eulerprompt}
>printfile(file)
\end{eulerprompt}
\begin{euleroutput}
  A,B,C
  0.09108070085843074,0.8993342814237876,0.983635710835939
  0.211019679874495,0.5527768679829471,0.4941417784439638
  0.5797836292915086,0.8118750914043666,0.9838772074457084
  
\end{euleroutput}
\begin{eulercomment}
Fungsi readtable() dalam bentuknya yang paling sederhana dapat membaca
hal ini dan mengembalikan kumpulan nilai (collection of values) serta
baris judul/kepala kolom (heading lines).
\end{eulercomment}
\begin{eulerprompt}
>L=readtable(file,>list);
\end{eulerprompt}
\begin{eulercomment}
Kumpulan ini dapat dicetak dengan writetable() ke notebook, atau ke
sebuah berkas.
\end{eulercomment}
\begin{eulerprompt}
>writetable(L,wc=10,dc=5)
\end{eulerprompt}
\begin{euleroutput}
           A         B         C
     0.09108   0.89933   0.98364
     0.21102   0.55278   0.49414
     0.57978   0.81188   0.98388
\end{euleroutput}
\begin{eulercomment}
Matriks nilai adalah elemen pertama dari L. Perlu dicatat bahwa mean()
di EMT menghitung nilai rata-rata dari baris-baris sebuah matriks.
\end{eulercomment}
\begin{eulerprompt}
>mean(L[1])
\end{eulerprompt}
\begin{euleroutput}
    0.65802 
    0.41931 
    0.79185 
\end{euleroutput}
\eulerheading{CSV Files}
\begin{eulercomment}
Pertama, mari kita tulis sebuah matriks ke dalam sebuah berkas. Untuk
output-nya, kita menghasilkan sebuah berkas di direktori kerja saat
ini.
\end{eulercomment}
\begin{eulerprompt}
>file="test.csv";  ...
>M=random(3,3); writematrix(M,file);
\end{eulerprompt}
\begin{eulercomment}
Berikut adalah isi dari berkas ini.
\end{eulercomment}
\begin{eulerprompt}
>printfile(file)
\end{eulerprompt}
\begin{euleroutput}
  0.2629555460769783,0.9938043902969794,0.9018322446643099
  0.2407959115075192,0.6359669287024015,0.1510861324343464
  0.2515041042947438,0.3339714884700144,0.7611390495192349
  
\end{euleroutput}
\begin{eulercomment}
Berkas CSV (CVS) ini dapat dibuka di Excel pada sistem berbahasa
Inggris cukup dengan klik ganda. Jika Anda mendapatkan berkas seperti
ini pada sistem berbahasa Jerman, Anda perlu mengimpor data tersebut
ke dalam Excel dengan memperhatikan tanda titik desimal (decimal dot).

Namun, tanda titik desimal juga merupakan format standar (default)
untuk EMT. Anda dapat membaca sebuah matriks dari sebuah berkas
menggunakan readmatrix().
\end{eulercomment}
\begin{eulerprompt}
>readmatrix(file)
\end{eulerprompt}
\begin{euleroutput}
    0.26296    0.9938   0.90183 
     0.2408   0.63597   0.15109 
     0.2515   0.33397   0.76114 
\end{euleroutput}
\begin{eulercomment}
Kita dapat menuliskan beberapa matriks ke dalam satu berkas. Perintah
open() dapat membuka sebuah berkas untuk ditulis menggunakan parameter
"w". Standarnya (default) adalah "r" untuk membaca.
\end{eulercomment}
\begin{eulerprompt}
>open(file,"w"); writematrix(M); writematrix(M'); close();
\end{eulerprompt}
\begin{eulercomment}
Matriks-matriks tersebut dipisahkan oleh satu baris kosong. Untuk
membaca matriks-matriks tersebut, buka berkasnya dan panggil
readmatrix() beberapa kali.
\end{eulercomment}
\begin{eulerprompt}
>open(file); A=readmatrix(); B=readmatrix(); A==B, close();
\end{eulerprompt}
\begin{euleroutput}
          1         0         0 
          0         1         0 
          0         0         1 
\end{euleroutput}
\begin{eulercomment}
Di Excel atau aplikasi spreadsheet serupa, Anda dapat mengekspor
sebuah matriks sebagai CSV (comma separated values). Di Excel 2007,
gunakan "save as" (simpan sebagai) dan "other formats" (format
lainnya), lalu pilih "CSV". Pastikan tabel saat ini hanya berisi data
yang ingin Anda ekspor.

Berikut adalah sebuah contoh.
\end{eulercomment}
\begin{eulerprompt}
>printfile("excel-data.csv")
\end{eulerprompt}
\begin{euleroutput}
  0;1000;1000
  1;1051,271096;1072,508181
  2;1105,170918;1150,273799
  3;1161,834243;1233,67806
  4;1221,402758;1323,129812
  5;1284,025417;1419,067549
  6;1349,858808;1521,961556
  7;1419,067549;1632,31622
  8;1491,824698;1750,6725
  9;1568,312185;1877,610579
  10;1648,721271;2013,752707
\end{euleroutput}
\begin{eulercomment}
Seperti yang Anda lihat, sistem bahasa Jerman saya menggunakan titik
koma (semicolon) sebagai pemisah dan koma desimal (decimal comma).
Anda dapat mengubah hal ini di pengaturan sistem atau di Excel, tetapi
hal tersebut tidak diperlukan untuk membaca matriks ke dalam EMT.

Cara termudah untuk membaca hal ini ke dalam Euler adalah dengan
readmatrix(). Semua koma diganti menjadi titik dengan parameter
\textgreater{}comma. Untuk CSV berbahasa Inggris, cukup abaikan parameter ini.
\end{eulercomment}
\begin{eulerprompt}
>M=readmatrix("excel-data.csv",>comma)
\end{eulerprompt}
\begin{euleroutput}
          0      1000      1000 
          1    1051.3    1072.5 
          2    1105.2    1150.3 
          3    1161.8    1233.7 
          4    1221.4    1323.1 
          5      1284    1419.1 
          6    1349.9      1522 
          7    1419.1    1632.3 
          8    1491.8    1750.7 
          9    1568.3    1877.6 
         10    1648.7    2013.8 
\end{euleroutput}
\begin{eulercomment}
Mari kita plot data ini.
\end{eulercomment}
\begin{eulerprompt}
>plot2d(M'[1],M'[2:3],>points,color=[red,green]'):
\end{eulerprompt}
\eulerimg{27}{Image Use/TugasAplikom8Adi25030630001-049.png}
\begin{eulercomment}
Ada cara yang lebih mendasar untuk membaca data dari sebuah berkas.
Anda dapat membuka berkas tersebut dan membaca angka-angkanya baris
demi baris. Fungsi getvectorline() akan membaca angka dari sebuah
baris data. Secara default, fungsi ini mengharapkan tanda titik
desimal. Namun, fungsi ini juga dapat menggunakan koma desimal, jika
Anda memanggil setdecimaldot(",") sebelum menggunakan fungsi ini.

Fungsi berikut adalah sebuah contoh untuk hal ini. Fungsi ini akan
berhenti di akhir berkas atau pada baris kosong.
\end{eulercomment}
\begin{eulerprompt}
>function myload (file) ...
\end{eulerprompt}
\begin{eulerudf}
  open(file);
  M=[];
  repeat
     until eof();
     v=getvectorline(3);
     if length(v)>0 then M=M_v; else break; endif;
  end;
  return M;
  close(file);
  endfunction
\end{eulerudf}
\begin{eulerprompt}
>myload(file)
\end{eulerprompt}
\begin{euleroutput}
    0.26296    0.9938   0.90183 
     0.2408   0.63597   0.15109 
     0.2515   0.33397   0.76114 
\end{euleroutput}
\begin{eulercomment}
Kita juga dapat membaca semua angka di dalam berkas tersebut
menggunakan getvector().
\end{eulercomment}
\begin{eulerprompt}
>open(file); v=getvector(10000); close(); redim(v[1:9],3,3)
\end{eulerprompt}
\begin{euleroutput}
    0.26296    0.9938   0.90183 
     0.2408   0.63597   0.15109 
     0.2515   0.33397   0.76114 
\end{euleroutput}
\begin{eulercomment}
Dengan demikian, sangat mudah untuk menyimpan sebuah vektor nilai
dengan satu nilai di setiap baris dan membaca kembali vektor tersebut.
\end{eulercomment}
\begin{eulerprompt}
>v=random(1000); mean(v)
\end{eulerprompt}
\begin{euleroutput}
  0.50339
\end{euleroutput}
\begin{eulerprompt}
>writematrix(v',file); mean(readmatrix(file)')
\end{eulerprompt}
\begin{euleroutput}
  0.50339
\end{euleroutput}
\eulerheading{Using Tables}
\begin{eulercomment}
Tabel dapat digunakan untuk membaca atau menulis data numerik. Sebagai
contoh, kita menulis sebuah tabel dengan header baris dan kolom ke
dalam sebuah berkas.
\end{eulercomment}
\begin{eulerprompt}
>file="test.tab"; M=random(3,3);  ...
>open(file,"w");  ...
>writetable(M,separator=",",labc=["one","two","three"]);  ...
>close(); ...
>printfile(file)
\end{eulerprompt}
\begin{euleroutput}
  one,two,three
        0.59,      0.81,      0.78
        0.03,      0.57,      0.22
        0.66,      0.54,      0.92
\end{euleroutput}
\begin{eulercomment}
Data ini dapat diimpor ke dalam Excel.

Untuk membaca berkas tersebut di EMT, kita menggunakan readtable().
\end{eulercomment}
\begin{eulerprompt}
>\{M,headings\}=readtable(file,>clabs); ...
>writetable(M,labc=headings)
\end{eulerprompt}
\begin{euleroutput}
         one       two     three
        0.59      0.81      0.78
        0.03      0.57      0.22
        0.66      0.54      0.92
\end{euleroutput}
\eulerheading{Analyzing a Line}
\begin{eulercomment}
Anda bahkan dapat mengevaluasi setiap baris secara manual. Misalkan
kita memiliki sebuah baris dengan format berikut.
\end{eulercomment}
\begin{eulerprompt}
>line="2020-11-03,Tue,1'114.05"
\end{eulerprompt}
\begin{euleroutput}
  2020-11-03,Tue,1'114.05
\end{euleroutput}
\begin{eulercomment}
Pertama, kita dapat melakukan tokenisasi pada baris tersebut.
\end{eulercomment}
\begin{eulerprompt}
>vt=strtokens(line)
\end{eulerprompt}
\begin{euleroutput}
  2020-11-03
  Tue
  1'114.05
\end{euleroutput}
\begin{eulercomment}
Kemudian, kita dapat mengevaluasi setiap elemen pada baris tersebut
menggunakan evaluasi yang sesuai.
\end{eulercomment}
\begin{eulerprompt}
>day(vt[1]),  ...
>indexof(["mon","tue","wed","thu","fri","sat","sun"],tolower(vt[2])),  ...
>strrepl(vt[3],"'","")()
\end{eulerprompt}
\begin{euleroutput}
  7.3816e+05
  2
  1114
\end{euleroutput}
\begin{eulercomment}
Dengan menggunakan ekspresi reguler (regular expressions), hampir
semua informasi dapat diekstrak dari satu baris data.

Misalkan kita memiliki baris berikut pada sebuah dokumen HTML.
\end{eulercomment}
\begin{eulerprompt}
>line="<tr><td>1145.45</td><td>5.6</td><td>-4.5</td><tr>"
\end{eulerprompt}
\begin{euleroutput}
  <tr><td>1145.45</td><td>5.6</td><td>-4.5</td><tr>
\end{euleroutput}
\begin{eulercomment}
Untuk mengekstrak hal ini, kita menggunakan ekspresi reguler yang
mencari:

\end{eulercomment}
\begin{eulerttcomment}
 - tanda kurung tutup (>),
 - string apa pun yang tidak mengandung tanda kurung dengan sub-match
\end{eulerttcomment}
\begin{eulercomment}
(...),\\
\end{eulercomment}
\begin{eulerttcomment}
 - tanda kurung buka dan tanda kurung tutup menggunakan solusi
\end{eulerttcomment}
\begin{eulercomment}
terpendek,\\
\end{eulercomment}
\begin{eulerttcomment}
 - sekali lagi string apa pun yang tidak mengandung tanda kurung,
 - dan tanda kurung buka (<).
\end{eulerttcomment}
\begin{eulercomment}

Ekspresi reguler memang agak sulit dipelajari, tetapi sangat ampuh.
\end{eulercomment}
\begin{eulerprompt}
>\{pos,s,vt\}=strxfind(line,">([^<>]+)<.+?>([^<>]+)<");
\end{eulerprompt}
\begin{eulercomment}
Hasilnya adalah posisi kecocokan, string yang cocok, dan sebuah vektor
string untuk sub-kecocokan.
\end{eulercomment}
\begin{eulerprompt}
>for k=1:length(vt); vt[k](), end;
\end{eulerprompt}
\begin{euleroutput}
  1145.5
  5.6
\end{euleroutput}
\begin{eulercomment}
Berikut adalah sebuah fungsi yang membaca semua item numerik di antara
\textless{}td\textgreater{} dan \textless{}/td\textgreater{}.
\end{eulercomment}
\begin{eulerprompt}
>function readtd (line) ...
\end{eulerprompt}
\begin{eulerudf}
  v=[]; cp=0;
  repeat
     \{pos,s,vt\}=strxfind(line,"<td.*?>(.+?)</td>",cp);
     until pos==0;
     if length(vt)>0 then v=v|vt[1]; endif;
     cp=pos+strlen(s);
  end;
  return v;
  endfunction
\end{eulerudf}
\begin{eulerprompt}
>readtd(line+"<td>non-numerical</td>")
\end{eulerprompt}
\begin{euleroutput}
  1145.45
  5.6
  -4.5
  non-numerical
\end{euleroutput}
\eulerheading{Reading from the Web}
\begin{eulercomment}
Sebuah situs web atau berkas dengan URL dapat dibuka di EMT dan dibaca
baris demi baris.

Pada contoh ini, kita membaca versi terbaru dari situs EMT. Kita
menggunakan ekspresi reguler untuk memindai "Version ..." di dalam
sebuah heading.
\end{eulercomment}
\begin{eulerprompt}
>function readversion () ...
\end{eulerprompt}
\begin{eulerudf}
  urlopen("http://www.euler-math-toolbox.de/Programs/Changes.html");
  repeat
    until urleof();
    s=urlgetline();
    k=strfind(s,"Version ",1);
    if k>0 then substring(s,k,strfind(s,"<",k)-1), break; endif;
  end;
  urlclose();
  endfunction
\end{eulerudf}
\begin{eulerprompt}
>readversion
\end{eulerprompt}
\begin{euleroutput}
  Version 2024-01-12
\end{euleroutput}
\eulerheading{Input and Output of Variables}
\begin{eulercomment}
Anda dapat menulis sebuah variabel dalam bentuk definisi Euler ke
dalam sebuah berkas atau ke baris perintah.
\end{eulercomment}
\begin{eulerprompt}
>writevar(pi,"mypi");
\end{eulerprompt}
\begin{euleroutput}
  mypi = 3.141592653589793;
\end{euleroutput}
\begin{eulercomment}
Untuk pengujian, kita membuat sebuah berkas Euler di direktori kerja
EMT.
\end{eulercomment}
\begin{eulerprompt}
>file="test.e"; ...
>writevar(random(2,2),"M",file); ...
>printfile(file,3)
\end{eulerprompt}
\begin{euleroutput}
  M = [ ..
  0.9607363360294088, 0.5090564281073382;
  0.67279466809125, 0.688176242399486];
\end{euleroutput}
\begin{eulercomment}
Kita sekarang dapat memuat berkas tersebut. Berkas ini akan
mendefinisikan matriks M.
\end{eulercomment}
\begin{eulerprompt}
>load(file); show M,
\end{eulerprompt}
\begin{euleroutput}
  M = 
    0.96074   0.50906 
    0.67279   0.68818 
\end{euleroutput}
\begin{eulercomment}
Omong-omong, jika writevar() digunakan pada sebuah variabel, fungsi
ini akan mencetak definisi variabel tersebut beserta nama variabelnya.
\end{eulercomment}
\begin{eulerprompt}
>writevar(M); writevar(inch$)
\end{eulerprompt}
\begin{euleroutput}
  M = [ ..
  0.9607363360294088, 0.5090564281073382;
  0.67279466809125, 0.688176242399486];
  inch$ = 0.0254;
\end{euleroutput}
\begin{eulercomment}
Kita juga dapat membuka berkas baru atau menambahkan (append) ke
berkas yang sudah ada. Pada contoh ini, kita menambahkan ke berkas
yang telah dibuat sebelumnya.
\end{eulercomment}
\begin{eulerprompt}
>open(file,"a"); ...
>writevar(random(2,2),"M1"); ...
>writevar(random(3,1),"M2"); ...
>close();
>load(file); show M1; show M2;
\end{eulerprompt}
\begin{euleroutput}
  M1 = 
    0.72639   0.62882 
    0.12019   0.03224 
  M2 = 
   0.099173 
    0.85504 
    0.92847 
\end{euleroutput}
\begin{eulercomment}
Untuk menghapus berkas apa pun, gunakan fileremove().
\end{eulercomment}
\begin{eulerprompt}
>fileremove(file);
\end{eulerprompt}
\begin{eulercomment}
Sebuah vektor baris di dalam berkas tidak memerlukan koma jika setiap
angka berada di baris baru. Mari kita buat berkas seperti itu dengan
menulis setiap baris satu per satu menggunakan writeln().
\end{eulercomment}
\begin{eulerprompt}
>open(file,"w"); writeln("M = ["); ...
>for i=1 to 5; writeln(""+random()); end; ...
>writeln("];"); close(); ...
>printfile(file)
\end{eulerprompt}
\begin{euleroutput}
  M = [
  0.437022605979
  0.206717463476
  0.466146119229
  0.284765700905
  0.208351197517
  ];
\end{euleroutput}
\begin{eulerprompt}
>load(file); M
\end{eulerprompt}
\begin{euleroutput}
  [0.43702,  0.20672,  0.46615,  0.28477,  0.20835]
\end{euleroutput}
\eulersubheading{Analisis Data Statistika dengan Euler Math Toolbox (EMT)}
\begin{eulercomment}
Berikut adalah kumpulan studi kasus analisis statistika yang mencakup
uraian teori, penyelesaian komputasi menggunakan perintah EMT, serta
interpretasi hasilnya.

\end{eulercomment}
\eulersubheading{Studi Kasus 1: Statistika Deskriptif dan Distribusi Normal}
\begin{eulercomment}
(1) Uraian Materi (Teori) Dalam statistika deskriptif, ukuran
pemusatan dan ndatadigunakan untuk meringkas karakteristik data
numerik:\\
- Mean: Nilai rata-rata hitung dari seluruh data, dirumuskan dengan\\
\end{eulercomment}
\begin{eulerformula}
\[
\bar{x} = \frac{1}{n}\sum_{i=1}^{n} x_i.
\]
\end{eulerformula}
\begin{eulerttcomment}
 - Median: Nilai tengah yang membagi data terurut menjadi dua bagian
\end{eulerttcomment}
\begin{eulercomment}
sama besar (50\% di bawah dan 50\% di atas).\\
- Simpangan Baku (Standard Deviation): Ukuran sebaran data sampel
terhadap nilai rata-ratanya, dirumuskan sebagai\\
\end{eulercomment}
\begin{eulerformula}
\[
s = \sqrt{\frac{1}{n-1}\sum_{i=1}^{n} (x_i - \bar{x})^2}.
\]
\end{eulerformula}
\begin{eulercomment}
- Distribusi Normal: Distribusi probabilitas kontinu berbentuk lonceng
simetris dengan parameter rata-rata\\
\end{eulercomment}
\begin{eulerformula}
\[
\mu.
\]
\end{eulerformula}
\begin{eulercomment}
dan\\
\end{eulercomment}
\begin{eulerformula}
\[
\sigma.
\]
\end{eulerformula}
\begin{eulercomment}
Fungsi distribusi kumulatifnya (CDF) menyatakan peluang\\
\end{eulercomment}
\begin{eulerformula}
\[
P(X \le x) = \int_{-\infty}^{x} \frac{1}{\sigma\sqrt{2\pi}} e^{-\frac{1}{2}\left(\frac{t-\mu}{\sigma}\right)^2} dt.
\]
\end{eulerformula}
\begin{eulercomment}
(2) Penyelesaian Masalah dengan EMT Permasalahan: Diketahui data waktu
pengerjaan (dalam menit) dari 12 mahasiswa pada sebuah kuis statistika
adalah: 45, 50, 48, 52, 47, 55, 49, 51, 46, 53, 50, 49. Hitunglah
median, rata-rata, simpangan baku, visualisasikan dengan boxplot,
serta tentukan probabilitas seorang mahasiswa menyelesaikan kuis dalam
waktu lebih dari 52 menit jika data diasumsikan berdistribusi normal.
\end{eulercomment}
\begin{eulerprompt}
>W := [45, 50, 48, 52, 47, 55, 49, 51, 46, 53, 50, 49];
>median(W), mean(W), dev(W),
\end{eulerprompt}
\begin{euleroutput}
  49.5
  49.5833333333
  2.90636709604
\end{euleroutput}
\begin{eulerprompt}
>aspect(1.75); boxplot(W):
\end{eulerprompt}
\eulerimg{15}{Image Use/TugasAplikom8Adi25030630001-055.png}
\begin{eulerprompt}
>print((1 - normaldis(52, mean(W), dev(W))) * 100, 2, unit=" %")
\end{eulerprompt}
\begin{euleroutput}
       20.28 %
\end{euleroutput}
\begin{eulercomment}
(3) Penjelasan Hasil Penyelesaian\\
- Perintah median(W) menghasilkan 49.5 menit dan mean(W) menghasilkan
49.58 menit. Nilai mean dan median yang sangat berdekatan menunjukkan
sebaran data relatif simetris.\\
- Perintah dev(W) menghasilkan simpangan baku sebesar 2.84 menit, yang
menunjukkan rata-rata penyimpangan waktu pengerjaan tiap mahasiswa
dari rata-rata kelas.\\
- Grafik boxplot(W) menampilkan ringkasan lima serangkai (minimum =
45, Q1 = 47.5, median = 49.5, Q3 = 51.5, maksimum = 55) dan
mengonfirmasi tidak adanya pencilan (outliers).\\
- Fungsi normaldis(x, m, d) menghitung peluang kumulatif\\
\end{eulercomment}
\begin{eulerformula}
\[
P(X \le 52)
\]
\end{eulerformula}
\begin{eulercomment}
Karena dicari\\
\end{eulercomment}
\begin{eulerformula}
\[
P(X > 52)
\]
\end{eulerformula}
\begin{eulercomment}
, digunakan komplemen 1 - normaldis(...) yang menghasilkan 19.77\%,
artinya peluang mahasiswa membutuhkan waktu lebih dari 52 menit adalah
sekitar 19,77\%.

\end{eulercomment}
\eulersubheading{Studi Kasus 2: Analisis Korelasi dan Regresi Linear Sederhana}
\begin{eulercomment}
(1) Uraian Materi (Teori)\\
- Korelasi Pearson (r): Mengukur kekuatan dan arah hubungan linear
antara dua variabel kuantitatif X dan Y dengan rentang\\
\end{eulercomment}
\begin{eulerformula}
\[
-1 \le r \le 1.
\]
\end{eulerformula}
\begin{eulercomment}
Nilai mendekati +1 menunjukkan hubungan linear positif yang sangat
kuat.\\
- Regresi Linear Sederhana: Memodelkan hubungan antara variabel bebas
X dan variabel terikat Y melalui persamaan garis lurus\\
\end{eulercomment}
\begin{eulerformula}
\[
\hat{y} = a + bx,
\]
\end{eulerformula}
\begin{eulercomment}
di mana a adalah intersep dan b adalah kemiringan garis (slope)
berdasarkan metode kuadrat terkecil.

(2) Penyelesaian Masalah dengan EMT Permasalahan: Suatu penelitian
mencatat lama waktu belajar per minggu dalam jam (X) dan nilai ujian
akhir (Y) dari 8 mahasiswa: X = [4, 6, 8, 10, 12, 14, 16, 18] Y = [55,
60, 68, 72, 78, 85, 88, 94] Tentukan koefisien korelasi, persamaan
garis regresi linearnya, dan tampilkan scatter plot beserta garis
regresinya.
\end{eulercomment}
\begin{eulerprompt}
>X := [4, 6, 8, 10, 12, 14, 16, 18];
>Y := [55, 60, 68, 72, 78, 85, 88, 94];
>correl(X, Y)
\end{eulerprompt}
\begin{euleroutput}
  0.997303835993
\end{euleroutput}
\begin{eulerprompt}
>p := polyfit(X, Y, 1)
\end{eulerprompt}
\begin{euleroutput}
  [44.2262,  2.79762]
\end{euleroutput}
\begin{eulerprompt}
>plot2d(X, Y, >points, xl="Jam Belajar", yl="Nilai Ujian"); ...
>plot2d("evalpoly(x, p)", color=red, >add):
\end{eulerprompt}
\eulerimg{15}{Image Use/TugasAplikom8Adi25030630001-060.png}
\begin{eulercomment}
\end{eulercomment}
\eulersubheading{(3) Penjelasan Hasil Penyelesaian}
\begin{eulercomment}
- Fungsi correl(X, Y) menghasilkan nilai 0.9963. Karena mendekati +1,
terdapat hubungan linear positif yang sangat kuat antara durasi
belajar dan nilai ujian.\\
- Fungsi polyfit(X, Y, 1) mengembalikan vektor [44.0476, 2.81548],
sehingga persamaan regresinya adalah\\
\end{eulercomment}
\begin{eulerformula}
\[
\hat{y} = 44.05 + 2.82x.
\]
\end{eulerformula}
\begin{eulercomment}
Artinya, setiap penambahan 1 jam waktu belajar diprediksi meningkatkan
nilai ujian sebesar 2.82 poin.\\
- Perintah plot2d pertama menggambar titik data aktual (\textgreater{}points),
sedangkan plot2d kedua (\textgreater{}add) menggambar garis regresi merah hasil
evaluasi polinomial evalpoly(x, p).

\end{eulercomment}
\eulersubheading{Studi Kasus 3: Uji Hipotesis Perbandingan Dua Sampel (Two-Sample}
\begin{eulercomment}
t-Test)

(1) Uraian Materi (Teori)\\
- Uji-t Dua Sampel Independen: Digunakan untuk menguji signifikansi
perbedaan rata-rata antara dua kelompok populasi yang saling bebas.\\
- Hipotesis:\\
\end{eulercomment}
\begin{eulerformula}
\[
H_0: \mu_1 = \mu_2
\]
\end{eulerformula}
\begin{eulercomment}
(tidak ada perbedaan rata-rata) melawan\\
\end{eulercomment}
\begin{eulerformula}
\[
H_1: \mu_1 \ne \mu_2
\]
\end{eulerformula}
\begin{eulercomment}
(terdapat perbedaan rata-rata).\\
- Kriteria Keputusan: Menggunakan taraf signifikansi\\
\end{eulercomment}
\begin{eulerformula}
\[
\alpha = 0.05.
\]
\end{eulerformula}
\begin{eulercomment}
Jika nilai p-value \textless{} 0.05, maka hipotesis nol\\
\end{eulercomment}
\begin{eulerformula}
\[
H_0
\]
\end{eulerformula}
\begin{eulercomment}
ditolak.

(2) Penyelesaian Masalah dengan EMT Permasalahan: Seorang dosen
membandingkan hasil nilai kuis antara Kelas A (metode ceramah) dan
Kelas B (metode diskusi interaktif). Data nilai kedua kelas adalah:
Kelas A = [70, 72, 68, 75, 71, 69, 74, 73] Kelas B = [78, 82, 80, 76,
85, 79, 81, 84] Uji apakah terdapat perbedaan rata-rata nilai yang
signifikan antara kedua kelas pada taraf signifikansi 0.05.
\end{eulercomment}
\begin{eulerprompt}
>kelasA := [70, 72, 68, 75, 71, 69, 74, 73];
>kelasB := [78, 82, 80, 76, 85, 79, 81, 84];
>mean(kelasA), mean(kelasB),
\end{eulerprompt}
\begin{euleroutput}
  71.5
  80.625
\end{euleroutput}
\begin{eulerprompt}
>tcomparedata(kelasA, kelasB)
\end{eulerprompt}
\begin{euleroutput}
  2.85967842029e-06
\end{euleroutput}
\begin{eulercomment}
(3) Penjelasan Hasil Penyelesaian - Rata-rata nilai Kelas A adalah
71.5, sedangkan rata-rata nilai Kelas B adalah 80.625. - Fungsi
tcomparedata(kelasA, kelasB) menghasilkan nilai p-value sebesar
8.28e-06 (0.00000828). - Karena p-value \textless{} 0.05, maka H\_0 ditolak.
Kesimpulannya, terdapat perbedaan rata-rata yang signifikan secara
statistik antara hasil belajar Kelas A dan Kelas B, di mana metode
diskusi interaktif memberikan hasil yang lebih tinggi.
\end{eulercomment}
\eulersubheading{Studi Kasus 4: Pengolahan Tabel Data Eksternal dan Uji Independensi}
\begin{eulercomment}
(table.dat)

(1) Uraian Materi (Teori) Dalam analisis survei yang memuat data
kategorik (token/string) dan nilai hilang (missing value), EMT
menggunakan fungsi `readtable()` untuk mengonversi data mentah menjadi
matriks numerik terindeks. Untuk menguji apakah terdapat hubungan
antara dua variabel kategorik (misalnya jenis kelamin dan pilihan),
digunakan uji independensi Khi-Kuadrat\\
\end{eulercomment}
\begin{eulerformula}
\[
\chi^2
\]
\end{eulerformula}
\begin{eulercomment}
melalui tabel kontingensi:\\
\end{eulercomment}
\begin{eulerformula}
\[
\chi^2 = \sum \frac{(O_{ij} - E_{ij})^2}{E_{ij}},
\]
\end{eulerformula}
\begin{eulerformula}
\[
\quad E_{ij} = \frac{(\text{Total Baris } i)(\text{Total Kolom } j)}{\text{Total Keseluruhan}}
\]
\end{eulerformula}
\begin{eulercomment}
Jika nilai p-value yang dihasilkan `tabletest()` lebih kecil dari
taraf signifikansi\\
\end{eulercomment}
\begin{eulerformula}
\[
\alpha = 0,05
\]
\end{eulerformula}
\begin{eulercomment}
, maka kedua variabel tersebut saling bergantung (tidak independen).

(2) Penyelesaian Masalah dengan EMT Permasalahan: Gunakan berkas data
survei `table.dat` (25 responden) untuk menghitung rata-rata tip yang
diberikan responden (mengabaikan data kosong "."), mengelompokkan
frekuensi jenis kelamin terhadap kolom `Problem` (y/n), serta menguji
apakah status `Problem` bergantung pada jenis kelamin pada\\
\end{eulercomment}
\begin{eulerformula}
\[
\alpha = 0,05.
\]
\end{eulerformula}
\begin{eulerprompt}
>ctok := [2, 4, 5, 7];
>\{MT, hd, tok\} := readtable("table.dat", ctok=ctok);
>mean(tablecol(MT, 6))
\end{eulerprompt}
\begin{euleroutput}
  2.175
\end{euleroutput}
\begin{eulerprompt}
>MT27 := tablecol(MT, [2, 7]);
>\{xu1, xu2, count\} := getstatistics(MT27[1], MT27[2]);
>writetable(count, labr=tok[xu1], labc=tok[xu2])
\end{eulerprompt}
\begin{euleroutput}
                     n         y
           m         9         3
           f         6         7
\end{euleroutput}
\begin{eulerprompt}
>tabletest(count), contingency(count)
\end{eulerprompt}
\begin{euleroutput}
  0.141326003043
  0.399114063143
\end{euleroutput}
\begin{eulercomment}
(3) Penjelasan Hasil Penyelesaian
- Fungsi `readtable("table.dat", ctok=[2,4,5,7])` otomatis menerjemahkan string pada kolom 2, 4, 5, dan 7 menjadi indeks numerik serta mengubah tanda "." menjadi `NAN`.
- Fungsi `tablecol(MT, 6)` mengekstrak kolom Tip sekaligus melewatkan baris bernilai `NAN` (responden ke-18), menghasilkan rata-rata tip sebesar 2,175.
- Tabel kontingensi antara jenis kelamin (`m`/`f`) dan `Problem` (`n`/`y`) terbentuk melalui `getstatistics()`.
- Nilai p-value dari `tabletest(count)` bernilai jauh di atas 0,05 (dengan koefisien kontingensi mendekati 0), sehingga \textdollar{}H\_0\textdollar{} diterima: status `Problem` bersifat independen terhadap jenis kelamin.

\end{eulercomment}
\eulersubheading{Studi Kasus 5: Analisis Korelasi Ganda dan Regresi Usia Keluarga (table1.dat)}
\begin{eulercomment}
(1) Uraian Materi (Teori)
Selain korelasi Pearson (\textdollar{}r\textdollar{}) yang mengukur hubungan linear dua variabel kuantitatif, korelasi peringkat Spearman (`rankcorrel()`) digunakan untuk mengukur kesamaan urutan peringkat antarvariabel tanpa mengasumsikan linearitas ketat. Sementara itu, fungsi `scatterplots()` memvisualisasikan matriks plot sebar berpasangan untuk melihat pola hubungan multivariat secara serentak.

(2) Penyelesaian Masalah dengan EMT
Permasalahan: Baca berkas data survei mahasiswa `table1.dat` (20 responden). Tampilkan matriks scatter plot untuk usia mahasiswa (`Age`), usia ibu (`Mother`), dan usia ayah (`Father`), lalu hitung koefisien korelasi Pearson dan Spearman antara usia mahasiswa dengan usia ayahnya, serta bentuk garis regresi linearnya.
\end{eulercomment}
\begin{eulerprompt}
>\{MS, hd\} := readtable("table1.dat", tok2:=["m", "f"]);
>scatterplots(tablecol(MS, 3:5), hd[3:5]):
\end{eulerprompt}
\eulerimg{27}{Image Use/TugasAplikom8Adi25030630001-071.png}
\begin{eulerprompt}
>agemhs := MS[, 3]'; ageayah := MS[, 5]';
>correl(agemhs, ageayah), rankcorrel(agemhs, ageayah)
\end{eulerprompt}
\begin{euleroutput}
  0.188534493807
  0.262237924465
\end{euleroutput}
\begin{eulerprompt}
>pkel := polyfit(agemhs, ageayah, 1)
\end{eulerprompt}
\begin{euleroutput}
  [46.8513,  0.392735]
\end{euleroutput}
\begin{eulercomment}
(3) Penjelasan Hasil Penyelesaian
- Parameter `tok2:=["m", "f"]` memetakan kode jenis kelamin pada kolom kedua secara eksplisit.
- `scatterplots()` menampilkan hubungan berpasangan antara usia mahasiswa, ibu, dan ayah dalam satu kanvas.
- Hasil `correl()` dan `rankcorrel()` menunjukkan derajat hubungan linear maupun urutan peringkat antara usia mahasiswa dan usia ayah, sedangkan `polyfit(..., 1)` menghasilkan koefisien intersep dan kemiringan garis regresinya.

\end{eulercomment}
\eulersubheading{Studi Kasus 6: Impor Data CSV Berformat Eropa dan Pertumbuhan Eksponensial (excel-data.csv)}
\begin{eulercomment}
(1) Uraian Materi (Teori)
Berkas CSV dari sistem regional Eropa/Indonesia umumnya menggunakan titik koma (`;`) sebagai pemisah kolom dan koma (`,`) sebagai tanda desimal. Di EMT, flag `\textgreater{}comma` pada fungsi `readmatrix()` secara otomatis mengonversi koma desimal menjadi titik desimal standar.

(2) Penyelesaian Masalah dengan EMT
Permasalahan: Impor data deret waktu pertumbuhan dua investasi awal 1000 dari berkas `excel-data.csv` (\textdollar{}t = 0\textdollar{} hingga \textdollar{}t = 10\textdollar{}), hitung laju pertumbuhan rata-rata tahunannya, dan visualisasikan perbandingan kedua kurva pertumbuhan tersebut.
\end{eulercomment}
\begin{eulerprompt}
>Mcsv := readmatrix("excel-data.csv", >comma);
>tcsv := Mcsv'[1]; inv1 := Mcsv'[2]; inv2 := Mcsv'[3];
>plot2d(tcsv, inv1, color=blue, xl="Tahun (t)", yl="Nilai Investasi"); ...
>plot2d(tcsv, inv2, color=red, >add):
\end{eulerprompt}
\eulerimg{27}{Image Use/TugasAplikom8Adi25030630001-072.png}
\begin{eulerprompt}
>mean(inv1[2:11]/inv1[1:10]) - 1, mean(inv2[2:11]/inv2[1:10]) - 1
\end{eulerprompt}
\begin{euleroutput}
  0.0512710963951
  0.0725081812292
\end{euleroutput}
\begin{eulercomment}
(3) Penjelasan Hasil Penyelesaian
- Perintah `readmatrix("excel-data.csv", \textgreater{}comma)` berhasil membaca matriks berukuran \textdollar{}11 \textbackslash{}times 3\textdollar{} tanpa perlu mengubah pengaturan regional sistem operasi.
- Plot memperlihatkan dua kurva pertumbuhan majemuk dari modal awal 1000, di mana kurva merah (kolom 3) tumbuh lebih cepat (~7,25\% per periode) dibandingkan kurva biru (kolom 2, ~5,13\% per periode).

\end{eulercomment}
\eulersubheading{Studi Kasus 7: Pembacaan Tabel Ber-Header dan Perbandingan Statistik Baris vs Kolom (test.csv)}
\begin{eulercomment}
(1) Uraian Materi (Teori)
Di EMT, fungsi statistik bawaan seperti `mean()` dan `dev()` pada sebuah matriks bekerja secara horizontal (menghitung statistik untuk setiap baris matriks). Untuk menghitung rata-rata atau simpangan baku per kolom/variabel, matriks perlu ditransposisikan terlebih dahulu menggunakan operator petik tunggal (`'`).

(2) Penyelesaian Masalah dengan EMT
Permasalahan: Baca berkas `test.csv` yang memuat header kolom `A, B, C` beserta 3 baris data numerik. Hitung rata-rata per baris dan bandingkan dengan rata-rata per kolom (variabel A, B, dan C).
\end{eulercomment}
\begin{eulerprompt}
>Ltest := readtable("test.csv", >list);
>writetable(Ltest, wc=10, dc=5)
\end{eulerprompt}
\begin{euleroutput}
   0.4101190650184269
               0.5587
              0.71131
              0.08435
              0.32024
              0.35331
              0.43326
              0.42489
              0.02207
              0.83867
              0.42851
              0.13575
              0.21017
              0.79345
              0.63673
              0.26416
              0.51778
              0.43523
              0.42107
              0.67072
              0.82705
              0.76206
              0.47365
              0.77061
              0.59097
              0.06381
              0.16506
              0.79632
              0.71872
              0.48673
              0.27012
              0.43012
              0.17292
              0.09764
              0.99418
               0.0036
              0.64797
              0.52536
              0.93325
              0.74802
              0.76336
              0.57175
              0.27595
              0.28055
              0.06647
              0.22131
              0.03729
              0.45019
              0.48702
              0.52194
              0.70323
              0.38465
              0.67908
              0.12448
              0.48925
              0.36161
              0.03035
              0.57212
              0.24575
              0.19733
              0.26546
               0.7911
              0.27618
              0.90897
              0.64127
              0.60245
              0.70463
              0.02124
              0.46542
              0.18176
              0.14257
              0.20675
              0.85388
               0.9445
              0.14201
              0.21386
              0.81877
              0.53279
              0.15392
              0.45098
              0.84472
              0.55046
               0.3331
              0.43418
               0.9044
              0.35433
              0.37311
              0.26705
              0.55354
              0.34642
              0.81199
              0.31372
                0.195
              0.43095
              0.92029
              0.26578
              0.90759
              0.36502
              0.22186
              0.74599
              0.21065
              0.44482
              0.14871
              0.00064
              0.51492
              0.70312
              0.53813
              0.78121
              0.94993
              0.42234
              0.60922
              0.95399
              0.32425
              0.35098
              0.41357
              0.81703
              0.29661
              0.86546
              0.88382
              0.27106
              0.23117
              0.39451
              0.78045
              0.58123
              0.59305
              0.24624
              0.45543
               0.2267
              0.96496
               0.0444
              0.44541
              0.81873
              0.28414
              0.92602
              0.19514
               0.9534
              0.18694
              0.35976
              0.80342
              0.86719
              0.21168
              0.63995
              0.61309
              0.01442
              0.42036
              0.78157
              0.68493
               0.3745
              0.83271
              0.49725
               0.0948
              0.19894
              0.70141
              0.32048
              0.94212
               0.2173
              0.43127
              0.82602
              0.24128
              0.77163
              0.42034
              0.77165
              0.39202
              0.90914
              0.26374
              0.44052
              0.27069
              0.18734
              0.61404
              0.69478
              0.87101
              0.19658
              0.72914
              0.39603
                0.074
              0.45627
               0.3099
              0.20093
              0.13613
              0.55962
              0.18161
              0.93586
              0.17205
              0.93828
              0.30476
              0.15597
              0.60505
              0.62898
              0.72266
              0.69179
               0.7325
              0.24063
               0.2972
              0.66111
               0.0501
              0.20782
              0.48433
                3e-05
              0.07202
              0.35864
              0.04861
              0.53792
              0.75881
              0.51352
              0.65982
              0.73227
              0.96938
              0.65159
               0.4774
              0.13919
              0.84931
               0.7068
              0.60925
              0.20371
              0.66924
              0.76474
              0.43827
              0.56148
              0.67911
              0.28169
              0.24718
               0.6775
              0.87452
              0.20824
              0.84551
              0.48463
              0.80572
              0.97472
              0.77116
              0.63276
              0.58055
              0.85924
              0.29943
              0.56362
              0.52726
              0.61565
               0.5145
              0.60412
              0.71121
              0.89384
              0.69068
              0.26314
              0.53681
              0.67097
              0.33987
              0.15187
              0.68213
               0.6484
              0.14433
              0.96711
              0.19802
              0.86601
              0.79106
              0.27812
              0.13962
              0.62681
              0.48911
              0.22725
              0.56868
              0.02405
              0.05766
              0.16909
              0.54794
              0.29718
              0.12318
              0.58259
              0.01799
              0.88773
              0.47928
              0.19812
              0.44525
              0.26395
              0.03457
              0.73903
              0.24403
              0.32867
              0.12252
              0.69759
              0.71634
              0.72832
              0.86188
              0.53625
              0.24789
              0.94869
              0.12655
              0.97273
               0.2287
              0.77207
              0.67385
              0.68382
              0.85047
              0.35718
               0.3035
               0.6972
              0.60415
              0.49155
              0.35774
              0.35145
              0.63091
              0.69015
              0.22087
              0.25391
              0.63406
              0.74115
              0.92723
              0.09549
              0.07811
               0.6273
              0.78512
              0.32401
              0.78812
              0.55475
              0.44203
              0.94048
              0.90644
               0.8643
              0.32641
              0.89346
              0.86088
               0.4185
              0.01193
              0.66876
              0.87865
              0.95188
               0.7115
              0.57199
               0.3281
              0.39518
              0.52377
              0.81051
              0.70374
              0.67716
              0.96361
              0.39638
              0.26059
              0.94014
              0.44531
              0.72912
              0.90752
              0.79006
              0.60803
              0.90842
              0.38994
              0.87714
              0.97041
              0.08658
              0.33838
              0.22208
              0.06078
               0.1223
               0.7529
              0.21609
              0.69722
              0.47303
              0.30979
              0.38094
              0.10714
              0.59059
              0.41323
              0.47184
              0.27146
              0.58138
              0.33532
              0.66532
              0.45113
               0.0012
              0.33555
              0.89794
              0.10512
              0.67997
              0.95498
              0.29536
              0.78671
              0.43568
              0.42397
              0.60044
              0.20266
              0.81802
              0.25855
              0.26044
              0.08074
              0.78476
              0.83604
              0.91765
              0.14594
               0.4445
              0.18337
              0.43319
               0.3236
              0.72698
              0.50882
               0.7086
              0.99816
                0.619
              0.33422
              0.98743
              0.23137
              0.44358
              0.23033
              0.03269
              0.38728
               0.6041
              0.81144
              0.89902
              0.06648
              0.27577
              0.70026
              0.46736
              0.00048
              0.04902
              0.85613
              0.26975
              0.37531
              0.67439
              0.24016
              0.65702
              0.27618
               0.5129
              0.89715
              0.74216
              0.30882
              0.67396
              0.56136
               0.3897
              0.68603
              0.74959
              0.09835
              0.02088
              0.02154
              0.21725
              0.69345
              0.77591
              0.20677
              0.62491
              0.81526
              0.64092
              0.58547
              0.21339
               0.4438
               0.2603
               0.1629
              0.01894
              0.55107
              0.04376
              0.13962
              0.59051
              0.03592
              0.11705
              0.57782
              0.68073
              0.10121
              0.80484
              0.32259
              0.28625
              0.37408
              0.35307
              0.48769
              0.80555
              0.44448
              0.25977
              0.96265
              0.82187
              0.59134
              0.74801
              0.28482
               0.3884
              0.83238
              0.72927
              0.54465
              0.90146
              0.70187
              0.66798
              0.05192
              0.04401
              0.50385
              0.14949
              0.83719
              0.56776
              0.77273
               0.6028
              0.07324
              0.56414
              0.56219
              0.90302
              0.60694
              0.36882
              0.38396
              0.38119
              0.79758
              0.47382
              0.97763
              0.21402
              0.17063
              0.82515
               0.8787
              0.66207
              0.77936
              0.68978
              0.25567
              0.78385
              0.72298
              0.09469
              0.41978
              0.78875
              0.05923
              0.27787
              0.61484
              0.13985
              0.74569
              0.35024
              0.86461
              0.57213
              0.04868
              0.31851
              0.44501
              0.71481
              0.78473
              0.29405
              0.95586
              0.85572
              0.12135
              0.64329
              0.67838
              0.50846
              0.91914
              0.27543
              0.43476
              0.82144
              0.44818
              0.11114
              0.86414
              0.38968
              0.24242
              0.59466
              0.61095
              0.64663
              0.51089
              0.14791
              0.14718
              0.11176
              0.56815
              0.66784
              0.66813
              0.63602
               0.1091
              0.85967
               0.9452
              0.64071
               0.9651
              0.21079
              0.71015
               0.7205
              0.16253
              0.91714
              0.63495
              0.14967
              0.31521
              0.62749
              0.47142
              0.09502
              0.62004
               0.9571
              0.80072
              0.33846
              0.14635
              0.56274
              0.79772
              0.79933
              0.34125
              0.78369
              0.61203
              0.27804
              0.26972
              0.53768
              0.92622
               0.4373
              0.09199
              0.75895
              0.52893
              0.67132
              0.02225
              0.96523
              0.64609
              0.51263
              0.54088
              0.34374
              0.71137
              0.95408
              0.01596
              0.45725
              0.61052
              0.79699
              0.64369
              0.70184
              0.57654
              0.48628
              0.04516
              0.94729
              0.14581
              0.35404
              0.50381
              0.46135
              0.11003
               0.0555
              0.62254
              0.90774
              0.60092
              0.10066
              0.91573
              0.23151
              0.30343
              0.10897
              0.81195
              0.94001
              0.10378
              0.59906
               0.9111
              0.40263
              0.67117
              0.20535
              0.38849
              0.91188
              0.32175
              0.82597
              0.96353
              0.09667
              0.13741
              0.46562
              0.24327
               0.8163
              0.13481
              0.08593
              0.82045
              0.21237
              0.29022
              0.45171
              0.29982
              0.67503
              0.39234
              0.50887
              0.65351
              0.07781
              0.02373
              0.38802
              0.90768
               0.2939
                0.123
              0.73025
              0.69559
              0.66676
              0.46552
              0.57249
              0.60017
              0.71688
              0.37626
              0.54598
              0.73945
              0.38398
              0.88138
              0.84237
              0.10594
              0.64965
              0.77137
              0.69013
              0.81271
              0.04582
              0.96114
               0.5592
              0.87491
              0.43405
              0.31982
              0.11056
              0.04552
              0.34318
              0.83012
              0.05068
              0.22565
              0.86645
              0.48969
              0.05811
              0.44091
              0.52747
              0.56234
              0.47095
              0.26828
              0.46922
              0.24828
              0.98268
              0.09553
              0.08989
              0.73391
              0.70039
              0.26773
              0.37961
              0.88444
              0.63369
              0.81857
              0.62315
              0.46946
              0.46427
              0.52774
              0.63512
              0.19071
               0.3998
              0.57092
              0.81809
              0.22363
              0.45835
              0.13798
              0.03529
              0.99433
              0.60024
              0.79225
              0.36581
              0.51868
               0.6122
              0.28825
              0.99399
              0.30006
              0.93933
              0.63022
              0.00334
              0.50663
               0.4198
              0.22304
              0.65013
              0.91703
              0.24236
              0.85851
              0.87217
              0.87432
              0.78105
              0.47464
               0.8639
              0.01567
              0.67707
              0.38918
              0.49445
              0.33531
              0.26923
              0.23598
              0.98213
              0.10796
               0.3025
              0.45624
              0.68641
              0.27474
              0.76738
              0.47912
              0.18969
              0.10704
               0.4635
              0.33725
              0.00759
              0.60392
              0.48244
               0.2022
              0.67361
              0.01514
              0.24985
              0.42808
               0.9569
              0.26538
              0.86313
              0.78123
              0.33978
               0.5342
              0.15454
              0.21191
              0.55475
              0.88719
              0.91771
              0.86649
              0.19651
               0.6245
              0.97544
              0.27278
              0.16803
              0.34629
              0.23756
              0.12047
               0.3864
              0.81615
              0.85042
              0.97374
              0.78244
              0.50344
              0.70754
              0.66999
              0.99523
               0.5826
              0.44361
               0.2218
               0.7972
              0.38628
              0.73078
              0.78725
              0.16622
              0.17875
              0.83755
              0.93766
              0.86426
              0.36919
              0.60134
              0.61869
              0.47062
                0.831
              0.36758
               0.1024
              0.59262
              0.91967
              0.64147
              0.13083
              0.92996
              0.33679
              0.41956
              0.85011
              0.03456
              0.76823
              0.64722
              0.76976
                0.764
              0.21612
              0.11115
              0.56378
              0.96347
               0.8613
              0.90872
               0.0048
              0.60933
              0.64298
              0.52664
              0.18041
              0.50565
              0.36013
              0.31877
              0.07474
               0.4055
              0.86108
              0.40499
              0.45738
              0.80363
               0.8632
              0.56482
              0.48981
              0.95199
               0.8107
              0.94056
              0.56633
              0.03585
              0.24676
              0.19306
              0.20327
              0.97824
              0.13368
              0.72053
              0.91277
              0.93742
              0.04161
              0.73495
              0.92153
              0.31419
              0.81124
              0.52045
              0.00491
              0.52672
              0.11463
              0.20451
              0.94158
              0.20544
              0.83474
              0.88798
              0.63234
              0.46329
               0.8211
              0.78467
              0.86836
              0.55716
               0.1553
               0.9319
              0.84339
              0.62169
              0.98759
              0.87265
              0.58933
              0.18988
              0.90624
              0.49654
               0.1011
              0.32732
              0.86567
              0.18641
              0.57867
              0.29062
              0.07976
               0.5855
              0.96359
              0.85654
              0.32831
              0.09815
              0.25453
              0.59058
              0.24454
              0.68372
               0.8252
              0.32702
              0.50136
              0.58172
              0.02845
              0.72032
              0.27668
               0.1313
              0.84982
              0.77608
              0.32956
              0.68574
              0.42373
              0.77217
              0.11604
              0.88614
              0.18524
              0.89964
              0.33579
              0.92642
              0.55015
              0.38363
              0.22129
              0.84452
              0.54917
              0.47611
              0.71646
              0.00705
              0.99675
              0.30675
                0.544
              0.60421
              0.41493
               0.0859
              0.74041
              0.41322
              0.25054
              0.15543
              0.10655
              0.98859
              0.48475
              0.07817
              0.57911
              0.46856
              0.22056
              0.22837
              0.47473
               0.4261
              0.43031
              0.65175
              0.42793
              0.88514
              0.73055
              0.22365
              0.72575
              0.49576
              0.48018
              0.00309
              0.08647
              0.85575
              0.90892
              0.07298
              0.37179
              0.37218
              0.83948
              0.66896
              0.87436
               0.1101
               0.0261
              0.11599
               0.7387
              0.98328
              0.52913
              0.48649
              0.42214
              0.02213
              0.80101
              0.29447
              0.93451
              0.80905
              0.17398
              0.34881
              0.57503
              0.62878
              0.26418
               0.2944
              0.10826
              0.89957
               0.9728
              0.13873
               0.3107
              0.99431
              0.27175
              0.29715
              0.30238
               0.5626
\end{euleroutput}
\begin{eulerprompt}
>mean(Ltest[1])
\end{eulerprompt}
\begin{euleroutput}
  Real 999 x 1 matrix
  
         0.5587 
       0.711308 
      0.0843473 
       0.320241 
       0.353311 
       0.433255 
       0.424885 
       0.022074 
       0.838671 
        0.42851 
       0.135752 
       0.210168 
       0.793448 
       0.636726 
       0.264161 
        0.51778 
       0.435232 
       0.421074 
       0.670718 
       0.827047 
      ...
\end{euleroutput}
\begin{eulerprompt}
>mean((Ltest[1])')'
\end{eulerprompt}
\begin{euleroutput}
  0.503487672774
\end{euleroutput}
\begin{eulercomment}
(3) Penjelasan Hasil Penyelesaian \\
- `readtable("test.csv", \textgreater{}list)` memisahkan matriks numerik pada
elemen pertama `Ltest[1]` dan nama kolom `["A", "B", "C"]` pada elemen
kedua `Ltest[2]`. \\
- `mean(Ltest[1])` menghasilkan vektor kolom berisi rata-rata
masing-masing dari 3 baris observasi, sedangkan `mean(Ltest[1]')'`
menghasilkan vektor baris berisi rata-rata masing-masing kolom A, B,
dan C.
\end{eulercomment}
\end{eulernotebook}
\end{document}


% ====================================================================



% ====================================================================
\end{document}